\documentclass[11pt,letterpaper]{article}
\usepackage[margin=1in]{geometry}
\usepackage[utf8]{inputenc}
\usepackage[T1]{fontenc}
\usepackage{amsmath,amssymb,amsthm,amsfonts}
\usepackage{mathtools}
\usepackage{hyperref}
\usepackage{enumitem}

\hypersetup{colorlinks=true,linkcolor=blue,citecolor=blue,urlcolor=blue}

\newtheorem{theorem}{Theorem}[section]
\newtheorem{proposition}[theorem]{Proposition}
\newtheorem{lemma}[theorem]{Lemma}
\newtheorem{corollary}[theorem]{Corollary}
\theoremstyle{definition}
\newtheorem{remark}[theorem]{Remark}

\newcommand{\bbR}{\mathbb{R}}
\newcommand{\bbC}{\mathbb{C}}
\newcommand{\bbZ}{\mathbb{Z}}
\newcommand{\bbH}{\mathbb{H}}
\newcommand{\cA}{\mathcal{A}}
\newcommand{\cG}{\mathcal{G}}
\newcommand{\cM}{\mathcal{M}}
\newcommand{\cH}{\mathcal{H}}
\newcommand{\fg}{\mathfrak{g}}
\newcommand{\Lap}{\Delta}
\newcommand{\dvol}{\,d\mathrm{vol}}

\title{Asymptotic product structure and collar essential spectrum on $\mathrm{SU}(2)$ instanton moduli}
\author{Vico Bonfioli\\\small\texttt{vicobonfioli@gmail.com}}
\date{\today}

\begin{document}
\maketitle

\begin{abstract}
On the moduli space $\cM_k(S^4_r)$ of charge-$k$ anti-self-dual $\mathrm{SU}(2)$ connections on the round $4$-sphere of radius $r$, equipped with the natural $L^2$ metric, we record two spectral results. First, the charge-one case: under the Habermann/Doi--Matsumoto--Matsumoto identification $\cM_1(S^4_r) \cong \bbH^5_r$, the $L^2$-Laplacian has purely continuous spectrum $[4/r^2, \infty)$, saturating McKean's inequality. Second, the higher-charge collar: a Schwinger-parametrized closed form for the BPST scale-derivative cross-term gives the exact metric asymptotic $48\pi^2 s_1 s_2/R^2 + O((s/R)^4)$, from which the codimension-$j$ Uhlenbeck collar of $\cM_k(S^4_r)$ has essential-spectrum bottom $4j/r^2$ via a Weyl quasi-mode argument on the asymptotic product of hyperbolic cusps.
\end{abstract}

\section{Introduction}

Let $\cM_k(S^4_r)$ denote the moduli space of charge-$k$ anti-self-dual $\mathrm{SU}(2)$ connections on the round $4$-sphere of radius $r$, equipped with the natural $L^2$ Hitchin-like metric inherited from the configuration space of connections (Groisser--Parker \cite{GP1987,GP1989}, Habermann \cite{Habermann1993}, Doi--Matsumoto--Matsumoto \cite{DMM1987}). The geometry of these moduli spaces is asymptotically hyperbolic near the Uhlenbeck strata: at $k=1$, the moduli space is globally a hyperbolic $5$-space; at $k \ge 2$, the geometry asymptotes to a Riemannian product of bubble factors and a residual moduli factor in each Uhlenbeck collar.

We make this picture quantitative in two steps. Theorem~\ref{thm:M1-spectrum} records the saturation of McKean's inequality on $\cM_1(S^4_r)$, with explicit radius dependence: $\lambda_0(\cM_1(S^4_r)) = 4/r^2$, equal to the bottom of continuous spectrum. Lemma~\ref{lem:cross-term} then gives an exact Schwinger-parametrized closed form for the $L^2$ inner product of two BPST scale-derivative tangent vectors on $\bbR^4$ centered at separation $R$ with scales $s_1, s_2$:
\[
\langle \partial A_1/\partial s_1,\;\partial A_2/\partial s_2 \rangle_{L^2(\bbR^4)} \;=\; 48\pi^2 s_1 s_2 \int_0^1 t(1-t)\,\frac{2X(t) - t(1-t)R^2}{X(t)^2}\,dt,
\]
with leading asymptotic $48\pi^2 s_1 s_2/R^2$ as $s_i \ll R$. The linear-in-$s_i$ vanishing controls the cross-term in the asymptotic-product collar geometry, and Theorem~\ref{thm:M_k-collar-conditional} extracts the codimension-$j$ collar essential-spectrum bottom $4j/r^2$ via a Weyl quasi-mode argument on the product of $j$ hyperbolic cusps.

\paragraph{What this paper does not do.} The collar essential bottom $4j/r^2$ is a statement about the spectrum localized to each Uhlenbeck stratum, not about the global $\lambda_0(\cM_k(S^4_r))$; by domain monotonicity the global bottom is bounded above by $4/r^2$ via the shallowest collar regardless of $k$ (Remark~\ref{rmk:not-global-bottom}). Determining the global $\lambda_0(\cM_k)$ for $k \ge 2$ is open. The body of the moduli space carries interior curvature variation that neither McKean nor the asymptotic-product argument bounds directly.

\paragraph{Computational verification.} The closed form \eqref{eq:lemma-cross-term-exact} is verified numerically against the original 4-dimensional integral by uniform-box Monte Carlo \cite{lemma-schwinger-script}.

\section{Setup}\label{sec:setup}

Let $(M, g)$ be a closed oriented Riemannian $4$-manifold, $G$ a compact simple Lie group, $P \to M$ a principal $G$-bundle. The Yang--Mills action $S_{YM}[A] = \tfrac{1}{2g_{YM}^2}\int_M |F_A|^2\dvol$ has linearized Hessian at any smooth ASD connection $A$ acting in Coulomb gauge ($d_A^* a = 0$) as the restriction of the covariant Laplacian to co-closed $1$-forms. The moduli space of charge-$k$ ASD connections $\cM_k(M) = \{A : F_A^+ = 0,\ c_2(P) = k\}/\cG$ inherits an $L^2$ Riemannian metric from the configuration space:
\[
g_{\cM_k}(a, b) \;=\; \int_M \langle a, b\rangle\dvol, \qquad a, b \in T_{[A]}\cM_k = \ker d_A^* \cap \ker d_A^+.
\]
The Uhlenbeck compactification adds boundary strata at which one or more instantons bubble off, with scale parameter $s_i \to 0$.

\section{The $k=1$ case: McKean saturation on $\cM_1(S^4_r)$}\label{sec:M1}

Let $S^4_r$ denote the round $4$-sphere of radius $r$.

\begin{proposition}[Habermann \cite{Habermann1993}, Doi--Matsumoto--Matsumoto \cite{DMM1987}]\label{prop:M1-iso}
The moduli space of charge-one $\mathrm{SU}(2)$ instantons on $S^4_r$ with the natural $L^2$ metric is isometric to a hyperbolic $5$-space of constant sectional curvature $-1/r^2$:
\[
\cM_1(S^4_r) \;\cong\; \bbH^5_r \;=\; \bigl(\mathrm{SO}(5,1)/\mathrm{SO}(5),\; r^2 g_{\mathrm{std}}\bigr),
\]
where $g_{\mathrm{std}}$ is the standard hyperbolic metric of curvature $-1$.
\end{proposition}

\noindent\emph{Normalization bridge.} We use the unreduced $L^2$ norm on $\mathfrak{su}(2)$-valued $1$-forms throughout. With this choice, the per-bubble scale-derivative self-norm of a BPST instanton on $\bbR^4$ is $\|\partial A/\partial s\|^2_{L^2(\bbR^4)} = 16\pi^2$ (used in Lemma~\ref{lem:cross-term} below), while the Habermann/DMM moduli metric of Proposition~\ref{prop:M1-iso} reads $g_{\cM_1} = (48\pi^2/\rho^2)(d\rho^2 + dx_0^2)$ in scale-position coordinates $(\rho, x_0) \in \bbR^+ \times \bbR^4 \subset \bbH^5$. The factor of $3$ between $16\pi^2$ and $48\pi^2$ reflects the $\eta$-tensor contraction $\sum_{a,\mu}\eta^a_{\mu\nu}\eta^a_{\mu\rho} = 3\delta_{\nu\rho}$ (equivalently, the trace over $\mathfrak{su}(2)$); curvature $-1/r^2$ then follows from the standard $r^2 g_{\mathrm{std}}$ rescaling. The two conventions are used consistently within their respective sections.

\begin{theorem}[Spectral bottom on $\cM_1(S^4_r)$]\label{thm:M1-spectrum}
The spectrum of the natural $L^2$-Laplacian on $\cM_1(S^4_r) \cong \bbH^5_r$ is purely continuous, equal to $[4/r^2, \infty)$, with no $L^2$ eigenfunctions. The spectral bottom is
\[
\lambda_0\bigl(\cM_1(S^4_r)\bigr) \;=\; \frac{4}{r^2} \;=\; \frac{(n-1)^2}{4}\cdot\frac{1}{r^2}
\]
with $n = 5$, saturating McKean's inequality \cite{McKean1970}.
\end{theorem}

\begin{proof}
On unit hyperbolic space $\bbH^5$ in upper-half-space coordinates $\{(x_1,\ldots,x_4, y) : y > 0\}$ with metric $g_{\mathrm{std}} = (dx^2 + dy^2)/y^2$, the volume density is $\sqrt{\det g} = y^{-5}$ and the Laplace--Beltrami operator on functions $f(y)$ depending only on the radial coordinate reduces, after the standard computation $\Lap_g f = (\det g)^{-1/2}\partial_y((\det g)^{1/2} g^{yy}\partial_y f)$, to
\[
\Lap_{\bbH^5_r} f \;=\; \tfrac{1}{r^2}\bigl[y^2 f''(y) - 3\,y\,f'(y)\bigr]\qquad\text{on radial $f$ on }\bbH^5_r.
\]
The radial generalized eigenfunctions $\phi_s(y) = y^s$ then satisfy $-\Lap_{\bbH^5}\phi_s = s(4-s)\phi_s$, maximized at $s = 2$ with eigenvalue $\mu = 4$. The volume measure $\dvol = y^{-5}\,dx\,dy$ gives $\int |\phi_s|^2\dvol = \int y^{2s-5}\,dx\,dy$, divergent for all $s$, so the bottom $\lambda_0(\bbH^5) = 4$ lies at the bottom of continuous spectrum. Rescaling the metric by $r^2$ rescales the Laplacian by $r^{-2}$, hence the spectrum by $r^{-2}$: $\lambda_0(\bbH^5_r) = 4/r^2$. McKean's inequality \cite{McKean1970} on an $n$-manifold with sectional curvature $K \le -\kappa^2$ gives $\lambda_0 \ge (n-1)^2\kappa^2/4$, equal to $4/r^2$ here; equality holds on constant-curvature spaces.
\end{proof}

\begin{remark}\label{rmk:radius-units}
The dimensional structure is explicit: $\lambda_0$ has dimensions $[\mathrm{length}]^{-2}$, scaling as $1/r^2$ with the radius of the underlying $S^4$. The dimensionless ratio is $\lambda_0 \cdot r^2 = 4$.
\end{remark}

\section{Higher charge: collar essential spectrum via the Schwinger cross-term lemma}\label{sec:collar}

For $k \ge 2$ the global $L^2$ metric on $\cM_k(S^4_r)$ is not known in closed form. The body of $\cM_k$ is not globally of constant curvature, so McKean's theorem does not directly apply globally. In each Uhlenbeck collar, however, the geometry asymptotes to a Riemannian product of bubble factors (each modelled by the cusp end of $\bbH^5_r$) and a residual background factor. The asymptotic product structure was identified qualitatively by Taubes \cite{Taubes1988} and developed in Donaldson--Kronheimer \cite{DK1990} \S 7.3; Lemma~\ref{lem:cross-term} below makes it quantitative for the $L^2$-metric scale-derivative cross-term.

\subsection{The Schwinger cross-term lemma}

\begin{lemma}[Cross-term decoupling in the Uhlenbeck collar; exact Schwinger closed form]\label{lem:cross-term}
Let $A_1$ and $A_2$ be two BPST $\mathrm{SU}(2)$ instantons on $\bbR^4$ centered at distinct points $x_1, x_2$ with scales $s_1, s_2$ and separation $R = |x_1 - x_2|$. Then the $L^2$ inner product of their scale-derivative tangent vectors admits the exact closed form
\begin{equation}\label{eq:lemma-cross-term-exact}
\langle \partial A_1/\partial s_1,\;\partial A_2/\partial s_2\rangle_{L^2(\bbR^4)}
\;=\; 48\pi^2 s_1 s_2 \int_0^1 t(1-t)\,\frac{2X(t) - t(1-t)R^2}{X(t)^2}\,dt,
\end{equation}
where
\begin{equation}\label{eq:lemma-X-def}
X(t) \;:=\; t(1-t)R^2 + (1-t)s_1^2 + t s_2^2.
\end{equation}
The integral evaluates in elementary closed form: with $\Sigma := R^2 + s_1^2 + s_2^2$ and
$D := \Sigma^2 - 4 s_1^2 s_2^2 = \bigl(R^2 + (s_1+s_2)^2\bigr)\bigl(R^2 + (s_1-s_2)^2\bigr)$,
\begin{equation}\label{eq:lemma-ct-elementary}
\langle \partial A_1/\partial s_1,\;\partial A_2/\partial s_2\rangle_{L^2(\bbR^4)}
\;=\; 48\pi^2 s_1 s_2\left[\frac{\Sigma}{D} \;-\; \frac{4 s_1^2 s_2^2}{D^{3/2}}\,\operatorname{arccosh}\!\Bigl(\frac{\Sigma}{2 s_1 s_2}\Bigr)\right].
\end{equation}
Moreover \eqref{eq:lemma-ct-elementary} is intrinsic to the hyperbolic geometry of Proposition~\ref{prop:M1-iso}: $\Sigma/(2 s_1 s_2) = \cosh d$, where $d$ is the hyperbolic distance between $(x_1, s_1)$ and $(x_2, s_2)$ in the upper-half-space model of $\bbH^5$ (unit curvature), so
\begin{equation}\label{eq:lemma-ct-hyperbolic}
\langle \partial A_1/\partial s_1,\;\partial A_2/\partial s_2\rangle_{L^2(\bbR^4)}
\;=\; 24\pi^2\;\frac{\cosh d\,\sinh d \;-\; d}{\sinh^3 d}\;=:\;24\pi^2\,G(d),
\end{equation}
a function of the hyperbolic distance alone. $G$ is smooth and positive on $(0, \infty)$, with $G(d) \to 2/3$ as $d \to 0^{+}$: the cross-term extends continuously to the coincidence limit with value $16\pi^2$, matching the diagonal norm $\|\partial A_1/\partial s_1\|^2_{L^2} = 16\pi^2$ below. As $d \to \infty$, $G(d) = 2e^{-d} + O(d\,e^{-3d})$, which recovers the leading term of \eqref{eq:lemma-asymptotic} via $e^{-d} = (1+o(1))\,s_1 s_2/R^2$ in the regime $s_1, s_2 \ll R$.
In the asymptotic regime $s_1, s_2 \ll R$, the leading behaviour with explicit next-order correction is
\begin{equation}\label{eq:lemma-asymptotic}
\langle \partial A_1/\partial s_1,\;\partial A_2/\partial s_2\rangle_{L^2(\bbR^4)}
\;=\; \frac{48\pi^2 s_1 s_2}{R^2}\left[1 \;-\; \frac{s_1^2 + s_2^2}{R^2}\right] \;+\; O\!\left(\frac{s_1 s_2 (s_1^2 + s_2^2)^2\,\bigl|\log(\max(s_1,s_2)/R)\bigr|}{R^6}\right).
\end{equation}
The next-order $-(s_1^2+s_2^2)/R^2$ coefficient is exactly $-1$. The remainder carries a $|\log(\max(s_i)/R)|$ factor at order $(s/R)^4$ via the symmetric $u^2 v^2$ cross-term in the dimensionless expansion (with $u = s_1/R$, $v = s_2/R$). The on-axis slice is rational: with $F(u,v) := I/(48\pi^2 s_1 s_2/R^2)$ for $s_1 s_2 > 0$, the limit $\lim_{v\to 0} F(u, v) = 1/(1+u^2)$ holds, with no logarithm in the $v\equiv 0$ slice. In particular, fixing $(R, s_2)$ and letting $s_1 \to 0$, the cross-term vanishes \emph{linearly} in $s_1$, and the corresponding $L^2$-metric component (after normalization by the diagonal norm $\|\partial A_1/\partial s_1\|^2_{L^2(\bbR^4)} = 16\pi^2$, which is $s_1$-independent) also vanishes linearly in $s_1$. \emph{Normalization note:} we use the unreduced $L^2$ norm on $\mathfrak{su}(2)$-valued $1$-forms (so the per-bubble $L^2$ self-norm is $16\pi^2$); the Habermann/Groisser--Parker normalization of the hyperbolic moduli metric (giving the curvature $-1/r^2$ via $g_{\cM_1} = (48\pi^2/\rho^2)(d\rho^2 + dx_0^2)$) differs by a fixed factor and is used consistently in Theorem~\ref{thm:M1-spectrum}.
\end{lemma}

\begin{proof}
After the $\eta$-tensor contraction identity $\sum_{a,\mu} \eta^a_{\mu\nu}\eta^a_{\mu\rho} = 3\delta_{\nu\rho}$ for the self-dual 't~Hooft symbols, the inner product reduces to
\[
\langle \partial A_1/\partial s_1,\;\partial A_2/\partial s_2\rangle_{L^2}
\;=\; 48 s_1 s_2 \int_{\bbR^4} \frac{x \cdot (x - R e_1)}{(|x|^2 + s_1^2)^2\,\bigl((x - R e_1)^2 + s_2^2\bigr)^2}\,d^4 x.
\]
Use the Schwinger parametrization $A^{-2} = \int_0^\infty \alpha\,e^{-\alpha A}\,d\alpha$ for each squared denominator. Set $\xi := x - \xi_*$ with the shift $\xi_* := (\alpha_2 R/(\alpha_1+\alpha_2))\,e_1$ (we rename the integration variable to $\xi$ to avoid clash with the upper-half-space coordinate $y$ used later in the spectral analysis). Completing the square in $x$ gives the integrand as $\exp(-(\alpha_1+\alpha_2)|\xi|^2 - \alpha_1\alpha_2 R^2/(\alpha_1+\alpha_2) - \alpha_1 s_1^2 - \alpha_2 s_2^2)$. Evaluate the Gaussian $\xi$-integrals using $\int_{\bbR^4} e^{-\sigma|\xi|^2}\,d^4\xi = \pi^2/\sigma^2$, $\int |\xi|^2 e^{-\sigma|\xi|^2}\,d^4\xi = 2\pi^2/\sigma^3$, $\int \xi_1 e^{-\sigma|\xi|^2}\,d^4\xi = 0$. Repartition $(\alpha_1, \alpha_2) \to (\sigma, t)$ with $\sigma = \alpha_1+\alpha_2$, $t = \alpha_2/\sigma$ (Jacobian $\sigma$), then integrate $\sigma$ from $0$ to $\infty$ using $\int_0^\infty e^{-\sigma X}\,d\sigma = 1/X$ and $\int_0^\infty \sigma e^{-\sigma X}\,d\sigma = 1/X^2$. The result is exactly \eqref{eq:lemma-cross-term-exact}.

For the elementary form \eqref{eq:lemma-ct-elementary}, write $a = s_1^2$, $b = s_2^2$, $c = R^2$ and $L(t) := a(1-t) + b\,t$, so that $X = c\,t(1-t) + L$ and the integrand of \eqref{eq:lemma-cross-term-exact} collapses pointwise:
\[
t(1-t)\,\frac{2X - t(1-t)c}{X^2} \;=\; \frac{(X - L)(X + L)}{c\,X^2} \;=\; \frac1c\Bigl(1 - \frac{L^2}{X^2}\Bigr).
\]
The quadratic $X(t) = -c\,t^2 + (c + b - a)\,t + a$ is positive on $[0,1]$ (both roots lie outside the interval, since $X(0) = a > 0$, $X(1) = b > 0$ and the leading coefficient is negative) with $(c+b-a)^2 + 4ac = \Sigma^2 - 4ab = D$. The standard antiderivatives give
\[
\int_0^1 \frac{dt}{X} = \frac{2}{\sqrt D}\,\operatorname{arccosh}\Bigl(\frac{\Sigma}{2\sqrt{ab}}\Bigr),\qquad
\int_0^1 \frac{X'\,dt}{X^2} = \frac1a - \frac1b,\qquad
\int_0^1 \frac{dt}{X^2} = \frac{(a-b)^2 + c(a+b)}{ab\,D} + \frac{2c}{D}\int_0^1 \frac{dt}{X}.
\]
Decomposing $L^2 = A\,X + B\,X' + C$ with rational coefficients ($A = -(b-a)^2/c$, and $B$, $C$ determined by matching the linear and constant terms), assembling the three integrals, and using the polynomial identity
\[
a^3 - a^2 b + 3a^2 c - a b^2 + 2abc + 3ac^2 + b^3 + 3b^2c + 3bc^2 + c^3 \;=\; \Sigma\,D,
\]
one obtains \eqref{eq:lemma-ct-elementary}. The substitution $\Sigma = 2 s_1 s_2 \cosh d$, $\sqrt D = 2 s_1 s_2 \sinh d$ then gives \eqref{eq:lemma-ct-hyperbolic}; that $d$ is the $\bbH^5$ distance is the upper-half-space formula $\cosh d = 1 + \bigl(|x_1 - x_2|^2 + (s_1 - s_2)^2\bigr)/(2 s_1 s_2) = \Sigma/(2 s_1 s_2)$. The limits of $G$ follow from $\cosh d \sinh d - d = \tfrac23 d^3 + O(d^5)$ and $\sinh^3 d = d^3 + O(d^5)$ at $d \to 0$, and from $\cosh d\sinh d - d = \tfrac14 e^{2d}(1 - 4d\,e^{-2d} + O(e^{-4d}))$, $\sinh^3 d = \tfrac18 e^{3d}(1 + O(e^{-2d}))$ at $d \to \infty$. The identities and \eqref{eq:lemma-ct-elementary}--\eqref{eq:lemma-ct-hyperbolic} are verified symbolically and to 50 significant digits against direct quadrature of \eqref{eq:lemma-cross-term-exact}, and to machine precision against the 4-D grid of \cite{lemma-ct-rust}, in \cite{lemma-ct-elementary}.

The asymptotic \eqref{eq:lemma-asymptotic} follows by expanding $1/X(t)$ in powers of $(s_i/R)^2$ in the integrand and integrating term by term against $t(1-t)$. The leading $1$ and the explicit $-(s_1^2+s_2^2)/R^2$ correction follow from elementary $\int_0^1 t(1-t)\,dt = 1/6$ identities; the $\log(s/R)$ at the next order arises from the endpoint behaviour of the $u^2 v^2$ contribution where $X(t) \sim t s_2^2$ as $t \to 0$ and $X(t) \sim (1-t) s_1^2$ as $t \to 1$, producing logarithmic divergence after the polynomial $t(1-t)$ damping is exhausted. Numerical verification of \eqref{eq:lemma-cross-term-exact} against the original 4-dimensional integral has been performed by adaptive Gauss--Kronrod cubature (after a reduction to two dimensions exploiting the $O(3)$ symmetry around the axis through $x_1, x_2$): the closed form agrees with the 4D integral to $\approx 2.5\,\mathrm{ulp}$ ($\approx 5.5\times 10^{-16}$ relative error) across a $5\times 5\times 4 + 6$ grid of $(s_1, s_2, R)$ values spanning four orders of magnitude in scale ratio \cite{lemma-ct-rust}. The on-axis closed form $F(u,0) = 1/(1+u^2)$ and the symmetric next-order coefficients $c_4 = d_4 = 8$ in $F(\epsilon, \epsilon) = 1 - 2\epsilon^2 + 8(1+\log\epsilon)\epsilon^4 + O(\epsilon^6\log\epsilon)$ are confirmed symbolically in \cite{lemma-ct-higher-order}.
\end{proof}

\subsection{Off-diagonal cross-blocks: position derivatives and the gauge correction}

To compare $g_{U_\epsilon^{(j)}}$ to a Riemannian product we must control \emph{all} bubble-bubble matrix elements of the $L^2$ metric, not only the scale-scale entries of Lemma~\ref{lem:cross-term}. The position-derivative tangent vectors require an additional ingredient: the \emph{horizontal} (Coulomb-gauge) projection. The bare partial derivative $\partial A/\partial x_0^\rho$ of a BPST instanton on $\bbR^4$ contains a long-range $|x-x_0|^{-2}$ tail (the rotation gauge field) which is \emph{not} square-integrable individually; only the horizontal projection onto $\ker d_A^*$ is. We record the closed forms and decay rates of the relevant bubble-bubble inner products in the bare and gauge-fixed normalizations.

\begin{lemma}[Off-diagonal bubble cross-block bounds]\label{lem:OD}
Let $A_1, A_2$ be two BPST $\mathrm{SU}(2)$ instantons on $\bbR^4$ centered at $x_1, x_2$ with separation $R = |x_1 - x_2|$ and scales $s_1, s_2$, and write $v_s^{(i)} = \partial A_i/\partial s_i$ and $v_{x,\rho}^{(i)} = \partial A_i/\partial x_i^\rho$ for the naive tangent vectors. Let $\hat{v}_{x,\rho}^{(i)}$ denote the $L^2$-orthogonal projection onto $\ker d_{A_i}^*$ (Coulomb-horizontal projection).
\begin{enumerate}
\item[(i)] \emph{(Scale-scale; Lemma~\ref{lem:cross-term}.)} $\langle v_s^{(1)}, v_s^{(2)}\rangle_{L^2} = 48\pi^2 s_1 s_2/R^2 + O(s_1 s_2 (s/R)^4 \log(s/R))$.
\item[(ii)] \emph{(Scale-position, bare.)} Only the axial component $\rho=1$ (along $x_2-x_1$) is nonzero by $O(3)$-symmetry; the orthogonal components vanish identically. The axial component admits an exact Schwinger representation (recorded in \cite{lemma-od-script}; involving $1/X(t)^2$-type denominators from the squared $1/D_i^2$ propagators after $\sigma$-integration). Its leading asymptotic in $s_1, s_2 \ll R$ is
\begin{equation}\label{eq:OD-bare-sx-asymp}
\langle v_s^{(1)}, v_{x,1}^{(2)}\rangle_{L^2(\bbR^4)} = \frac{72\pi^2 s_1}{R} - 4\pi^2 s_1 R + O(s_1\,s_2^2/R).
\end{equation}
The linear-in-$R$ piece is the long-range tail of $v_{x,1}^{(2)}$ before gauge fixing. The exact closed form is not used in the sequel; only the gauge-fixed bound (iv) below enters the proofs of Theorems~\ref{thm:M_k-collar-conditional}--\ref{thm:M_k-collar-spectral-type}.
\item[(iii)] \emph{(Position-position, bare.)} The naive inner product $\langle v_{x,\rho}^{(1)}, v_{x,\sigma}^{(2)}\rangle_{L^2(\bbR^4)}$ is logarithmically infrared-divergent for $s_1, s_2 > 0$: the integrand contains a $\int (x\cdot(x-Re_1))^2/D_1^2 D_2^2\,d^4x$ piece with $|x|^{-4}$ tail at infinity. The horizontal projections $\hat{v}_{x,\rho}^{(i)}$, however, are square-integrable.
\item[(iv)] \emph{(Gauge-fixed off-diagonal bound.)} In the moduli $L^2$-metric, the gauge-fixed bubble-bubble cross-block has matrix entries bounded by
\begin{equation}\label{eq:OD-horiz}
|\langle \hat{v}_\alpha^{(1)}, \hat{v}_\beta^{(2)}\rangle_{L^2}| \le \frac{C\, s_1 s_2}{R^2}\,\|\hat{v}_\alpha^{(1)}\|_{L^2}\cdot\|\hat{v}_\beta^{(2)}\|_{L^2}\cdot\bigl(1 + |\log(\min(s_1,s_2)/R)|\bigr),
\end{equation}
uniformly in $s_1, s_2 \le c R$ for some absolute $C, c > 0$, with $\alpha, \beta \in \{s, x_1, x_2, x_3, x_4\}$ indexing the five bubble tangent directions and $\hat{v}_s = v_s$.

\item[(v)] \emph{(Bubble--background cross-block.)} Let $a$ be any tangent vector to the residual moduli factor $\cM_{k-j}$, supported (in $S^4$ coordinates) at distance $\ge R$ from a bubble center $x_i$ with bubble scale $s_i$. Then $|\langle a, \hat{v}_\alpha^{(i)}\rangle_{L^2}| \le C\,(s_i^2/R)\,\|a\|_{L^2}\cdot\|\hat{v}_\alpha^{(i)}\|_{L^2}$.
\end{enumerate}
\end{lemma}

\begin{proof}
Part (i) is Lemma~\ref{lem:cross-term}. For (ii), the BPST derivative computation gives
\[
v_s^{(1)} = -4s_1 \eta^a_{\mu\nu} x^\nu/D_1^2,\quad v_{x,\rho}^{(2)} = -2\eta^a_{\mu\rho}/D_2 + 4\eta^a_{\mu\nu}(x-Re_1)^\nu(x-Re_1)^\rho/D_2^2,
\]
with $D_i = |x-x_i|^2 + s_i^2$. The 't Hooft contraction $\sum_{a,\mu}\eta^a_{\mu\nu}\eta^a_{\mu\rho} = 3\delta_{\nu\rho}$ collapses the spin index and pins the surviving direction to $\rho=1$ (the line $x_1 x_2$); the transverse $\rho\in\{2,3,4\}$ integrands are odd under reflection $x_\rho \to -x_\rho$ and vanish. The two resulting scalar integrals
\[
24 s_1\!\int_{\bbR^4}\!\tfrac{x_1}{D_1^2 D_2}\,d^4x,\qquad -48 s_1\!\int_{\bbR^4}\!\tfrac{(x\cdot(x-Re_1))(x_1-R)}{D_1^2 D_2^2}\,d^4x
\]
are evaluated by the Schwinger parametrization of Lemma~\ref{lem:cross-term} (each squared denominator $1/D^2 = \int_0^\infty \alpha e^{-\alpha D}d\alpha$, each linear denominator $1/D = \int_0^\infty e^{-\alpha D}d\alpha$). Completing the square in $x$ with $\sigma = \alpha_1+\alpha_2$, $t = \alpha_2/\sigma$ shifts the Gaussian center to $\xi_* = (\alpha_2 R/\sigma)e_1$, and the standard Gaussian moments $\int e^{-\sigma|\xi|^2} = \pi^2/\sigma^2$, $\int |\xi|^2 e^{-\sigma|\xi|^2} = 2\pi^2/\sigma^3$, $\int \xi_1^2 e^{-\sigma|\xi|^2} = \pi^2/(2\sigma^3)$ produce an exact closed form involving denominators $X(t)^2$ (recorded in \cite{lemma-od-script}). The leading asymptotic \eqref{eq:OD-bare-sx-asymp} follows by substituting $X(t) \to t(1-t)R^2$ and using $\int_0^1 t(1-t)\,dt = 1/6$, $\int_0^1 (1-t)\,dt = 1/2$, $\int_0^1 t(1-t)^2\,dt = 1/12$. Numerical verification \cite{lemma-od-script} confirms \eqref{eq:OD-bare-sx-asymp} to $<10^{-3}$ relative error across $(R, s_1, s_2)$ ranging over four orders of magnitude.

For (iii), $|v_{x,\rho}^{(i)}|^2 \sim 4|\eta_{\mu\rho}^a|^2/D_i^2 \sim 12/|x|^4$ at large $|x|$, so $\int |v_{x,\rho}|^2\,d^4x$ diverges logarithmically at infinity. The trace integrand of $\langle v_{x,\rho}^{(1)}, v_{x,\sigma}^{(2)}\rangle$ inherits a $|x|^{-4}$ infrared tail directly.

For (iv), the horizontal projection $\hat{v} = v - d_{A} G_A d_A^* v$, where $G_A$ is the Green's function for the Laplacian on $\fg$-valued functions in the BPST background, kills the long-range $|x-x_0|^{-2}$ piece of $v_{x,\rho}$, leaving $|\hat{v}_{x,\rho}(x)| = O(s^2/|x-x_0|^3)$ at infinity (the field-strength rate). The cross-block bound \eqref{eq:OD-horiz} then follows from a weighted Sobolev pairing argument in the Taubes/Donaldson--Kronheimer framework, with the $\log$ factor entering through a borderline radial integral over the inter-bubble annulus $\min(s_1,s_2) \lesssim |x-x_i| \lesssim R$. The full proof, including the scale-uniformity of the weighted Coulomb estimate and the tracking of the $\log$ factor, is given in Appendix~\ref{app:weighted-sobolev}.

For (v): we apply the same horizontal-projection localization (Appendix~\ref{app:weighted-sobolev}, \S\ref{app:sec:bg}). The tangent vector $\hat{v}_\alpha^{(i)}$ has $L^2$ mass concentrated in a ball of radius $O(s_i)$; pairing with a residual tangent vector $a$ supported at distance $\ge R$ from $x_i$ gives the stated $s_i^2/R$ rate via the four-dimensional Sobolev embedding $H^2 \hookrightarrow L^\infty$ together with the Coulomb-tail estimate of step (iv).
\end{proof}

\subsection{Collar essential-spectrum bottom}

\begin{theorem}[Collar-\emph{localized} essential-spectrum bottom on $\cM_k(S^4_r)$]\label{thm:M_k-collar-conditional}
Let $\cM_k(S^4_r)$ be the $k$-instanton moduli for $\mathrm{SU}(2)$ on $S^4_r$, equipped with the $L^2$ metric, and let $U_\epsilon^{(j)} \subset \cM_k$ denote the tubular neighborhood of the Uhlenbeck stratum at which exactly $j$ instantons have scale parameter below $\epsilon$ (the codimension-$j$ collar). Then the collar-\emph{localized} essential-spectrum bottom satisfies
\[
\lim_{\epsilon \to 0}\lambda_0^{\mathrm{ess}}\bigl(U_\epsilon^{(j)}\bigr) \;=\; \frac{4j}{r^2}\qquad\text{for each } 1 \le j \le k.
\]
In particular, the shallowest collar (one bubbling instanton) contributes essential spectrum at $4/r^2$, and the deepest stratum (all $k$ instantons simultaneously bubbling) contributes essential spectrum at $4k/r^2$.
\end{theorem}

\begin{remark}[Domain monotonicity bounds the global bottom]\label{rmk:not-global-bottom}
By domain monotonicity, $\lambda_0(\cM_k) \le \lambda_0^{\mathrm{ess}}(U_\epsilon^{(1)}) \to 4/r^2$, so $\lambda_0(\cM_k(S^4_r)) \le 4/r^2$ for every $k \ge 1$. Theorem~\ref{thm:M_k-collar-conditional} identifies the essential-spectrum contributions of each Uhlenbeck stratum but does \emph{not} settle whether $\lambda_0(\cM_k) = 4/r^2$ for all $k$, nor whether interior $L^2$ eigenvalues below this threshold exist. The bound $4k/r^2$ at the deepest stratum is the essential-spectrum bottom of a deep slice, not the bottom of the full Laplacian; the latter is bounded above by $4/r^2$ regardless. The two are distinct quantities and should not be conflated.
\end{remark}

\begin{proof}[Proof of Theorem~\ref{thm:M_k-collar-conditional}]
By Lemma~\ref{lem:OD} (clauses (i), (iv), (v)) the $L^2$ metric on $U_\epsilon^{(j)}$ admits an off-diagonal cross-block bound in the gauge-fixed (horizontal) basis of bubble tangent directions: for each pair of bubble-tangent matrix entries (scale-scale, scale-position, position-position) and each bubble-background matrix entry, the absolute value is bounded by a constant multiple of $\epsilon\,|\log\epsilon|$ times the geometric mean of the diagonal entries on either side, where $\epsilon = \max_i s_i$. Consequently, in any orthonormal basis adapted to the diagonal block structure (bubble $i$ and background $\cM_{k-j}$), the metric tensor satisfies the operator-norm comparison
\begin{equation}\label{eq:metric-product-bound}
\bigl(1 - C\epsilon|\log\epsilon|\bigr)\Bigl(g_{\cM_{k-j}(S^4_r)} \oplus \bigoplus_{i=1}^j g_{\mathrm{bubble},i}\Bigr) \;\le\; g_{U_\epsilon^{(j)}} \;\le\; \bigl(1 + C\epsilon|\log\epsilon|\bigr)\Bigl(g_{\cM_{k-j}(S^4_r)} \oplus \bigoplus_{i=1}^j g_{\mathrm{bubble},i}\Bigr),
\end{equation}
for some constant $C$ depending only on the geometry of $S^4_r$ (in particular on the lower bound for the separation $R$ in the collar, which can be taken bounded away from zero throughout $U_\epsilon^{(j)}$ for $\epsilon$ small enough). By Proposition~\ref{prop:M1-iso}, each individual bubble factor is, in the limit $s_i \to 0$, isometric to the cusp end (horosphere bundle near infinity) of $\bbH^5_r$.

In fact \eqref{eq:metric-product-bound} extends to a $C^2$ comparison of the perturbation tensor itself. Writing $h := g_{U_\epsilon^{(j)}} - g_{\mathrm{prod}}$ and denoting by $\nabla^{\mathrm{prod}}$ the Levi-Civita connection of the product metric, for $\epsilon$ small enough one has, uniformly on the cusp ends of $U_\epsilon^{(j)}$:
\begin{equation}\label{eq:metric-product-bound-C2}
|h|_{g_{\mathrm{prod}}} \;\le\; C\epsilon|\log\epsilon|, \qquad
|\nabla^{\mathrm{prod}} h|_{g_{\mathrm{prod}}} \;\le\; C\,\frac{\epsilon|\log\epsilon|}{r_{\mathrm{cusp}}}, \qquad
|(\nabla^{\mathrm{prod}})^{2} h|_{g_{\mathrm{prod}}} \;\le\; C\,\frac{\epsilon|\log\epsilon|^{2}}{r_{\mathrm{cusp}}^{2}},
\end{equation}
where $r_{\mathrm{cusp}}$ is the (inverse-length) cusp scale on the relevant bubble factor (in upper-half-space coordinates, $r_{\mathrm{cusp}}^{-1} \sim y_i\partial_{y_i}$ acting as a $C^0$-bounded unit). This $C^2$ statement is the $n = 0, 1, 2$ specialization of Lemma~\ref{app:lem:iterated-h}: each application of $\nabla^{\mathrm{prod}}$ in cusp-adapted coordinates is equivalent (up to bounded coefficients of the $\bbH^5_r$ connection symbols) to a single application of the cusp-dilation derivative $y_i\partial_{y_i}$, which by Lemma~\ref{app:lem:iterated-h} costs at most one extra logarithmic factor per derivative and no inverse power of the bubble scale.

\emph{Upper bound on the essential bottom: Weyl quasi-modes.}
The cusps of $\bbH^5_r$ have purely continuous spectrum with no $L^2$ ground state. In upper-half-space coordinates $(\xi^{(i)}_1, \ldots, \xi^{(i)}_4, y_i)$ for the $i$-th cusp ($i = 1, \ldots, j$), with $y_i > 0$, the radial generalized eigenfunctions of the Laplace--Beltrami operator on each cusp are $\phi^{(i)}_s(y_i) = y_i^s$, with eigenvalue $r^{-2}s(4-s)$ maximized at $s = 2$ with value $4/r^2$. Choose smooth cutoffs $\chi_n \in C_c^\infty(\bbR)$ supported in $[-n, n]$ with $\chi_n \equiv 1$ on $[-n+1, n-1]$ and $\|\chi_n'\|_\infty + \|\chi_n''\|_\infty \le C$. Define
\[
\Phi^{(j)}_n(y_1, \ldots, y_j) \;=\; \prod_{i=1}^j y_i^2 \chi_n(\log y_i).
\]
Direct calculation of $-\Lap\Phi^{(j)}_n$ on the \emph{exact} product cusp gives $(-\Lap_{\mathrm{prod}} + 4j/r^2)\Phi^{(j)}_n = O(\chi_n') + O(\chi_n'')$, and the routine estimate $\|\chi_n\|_{L^2} \sim \sqrt{n}$, $\|\chi_n'\|_{L^2} = O(1)$ yields
\[
\frac{\|(-\Lap_{\mathrm{prod}} - 4j/r^2)\Phi^{(j)}_n\|_{L^2}}{\|\Phi^{(j)}_n\|_{L^2}} \;\longrightarrow\; 0\qquad\text{as } n \to \infty.
\]
We transport this quasi-mode into $U_\epsilon^{(j)}$ using \eqref{eq:metric-product-bound}. The Laplace--Beltrami operator $\Lap_{U_\epsilon^{(j)}}$ on the collar differs from $\Lap_{\mathrm{prod}}$ by an operator of the form $L_\epsilon = (g_{U_\epsilon}^{-1} - g_{\mathrm{prod}}^{-1})\partial\partial + \text{lower-order}$, with $\|g_{U_\epsilon}^{-1} - g_{\mathrm{prod}}^{-1}\|_{\mathrm{op}} \le C'\epsilon|\log\epsilon|$ by \eqref{eq:metric-product-bound}. Hence
\[
\frac{\|(-\Lap_{U_\epsilon^{(j)}} - 4j/r^2)\Phi^{(j)}_n\|_{L^2}}{\|\Phi^{(j)}_n\|_{L^2}} \;\le\; \frac{\|(-\Lap_{\mathrm{prod}} - 4j/r^2)\Phi^{(j)}_n\|_{L^2}}{\|\Phi^{(j)}_n\|_{L^2}} + C'\epsilon|\log\epsilon|\cdot\frac{\|\nabla^2 \Phi^{(j)}_n\|_{L^2}}{\|\Phi^{(j)}_n\|_{L^2}}.
\]
The second term involves $\|\nabla^2\Phi^{(j)}_n\|_{L^2}^2/\|\Phi^{(j)}_n\|_{L^2}^2 = O(1)$ as $n\to\infty$. Explicitly, the hyperbolic Laplacian on the cusp acts on $y_i^2 = e^{2\log y_i}$ by an order-one operator in the logarithmic coordinate $u_i = \log y_i$ (since $y_i\partial_{y_i} = \partial_{u_i}$ and the metric is $dx_i^2/y_i^2 + du_i^2$ along the cusp), so $\|\nabla^2(y_i^2\chi_n(u_i))\|_{L^2}^2 = O(\|\chi_n\|_{L^2}^2 + \|\chi_n'\|_{L^2}^2 + \|\chi_n''\|_{L^2}^2) = O(n + 1)$, while $\|y_i^2\chi_n(u_i)\|_{L^2}^2 = O(n)$ (the $y_i^4 \cdot y_i^{-5}$ factor integrates against $du_i\,dx_i$ to give $O(\|\chi_n\|_{L^2}^2)$). The ratio is therefore $O(1)$, with the $\chi_n', \chi_n''$ contributions producing the $o(1)$ part of the first term in the displayed bound. Choosing $n = n(\epsilon)$ with $n(\epsilon) \to \infty$ slowly enough that the first term is $o(1)$ while $\epsilon|\log\epsilon|\to 0$ ensures both terms vanish. Hence $4j/r^2 \in \sigma_{\mathrm{ess}}(-\Lap_{U_\epsilon^{(j)}})$ for $\epsilon$ small enough, giving the upper bound $\liminf_\epsilon \lambda_0^{\mathrm{ess}}(U_\epsilon^{(j)}) \le 4j/r^2$.

\emph{Lower bound: form comparison and min-max.}
For the matching lower bound, the operator inequality \eqref{eq:metric-product-bound} on the metric tensor implies the comparison of Laplace--Beltrami operators as quadratic forms on $C_c^\infty(U_\epsilon^{(j)})$:
\[
\bigl(1 - C\epsilon|\log\epsilon|\bigr)\langle -\Lap_{\mathrm{prod}}\,\phi, \phi\rangle \le \langle -\Lap_{U_\epsilon^{(j)}}\,\phi, \phi\rangle \le \bigl(1 + C\epsilon|\log\epsilon|\bigr)\langle -\Lap_{\mathrm{prod}}\,\phi, \phi\rangle.
\]
By the min-max principle, $\lambda_0^{\mathrm{ess}}(-\Lap_{U_\epsilon^{(j)}}) \ge (1 - C\epsilon|\log\epsilon|)\lambda_0^{\mathrm{ess}}(-\Lap_{\mathrm{prod}})$. The product Laplacian decomposes as $-\Lap_{\mathrm{prod}} = -\Lap_{\cM_{k-j}} \otimes 1 + 1\otimes(-\Lap_{\mathrm{bubbles}})$, and the spectrum of a tensor sum is the Minkowski sum of the constituent spectra; thus $\lambda_0^{\mathrm{ess}}(-\Lap_{\mathrm{prod}}) \ge \lambda_0(\cM_{k-j}) + \sum_{i=1}^j \lambda_0(\bbH^5_r\text{-cusp}) \ge 0 + 4j/r^2$, using $\lambda_0 \ge 0$ on the (possibly compact) residual factor and the McKean saturation $\lambda_0(\bbH^5_r) = 4/r^2$ from Theorem~\ref{thm:M1-spectrum}. Taking $\epsilon \to 0$ gives $\liminf_\epsilon \lambda_0^{\mathrm{ess}}(U_\epsilon^{(j)}) \ge 4j/r^2$. The two bounds match.
\end{proof}

\begin{corollary}[Quantitative two-sided rate]\label{cor:rate}
There exist $C = C(j, r) > 0$ and $\epsilon_0 > 0$ such that for all $0 < \epsilon < \epsilon_0$,
\[
\Bigl|\lambda_0^{\mathrm{ess}}\bigl(U_\epsilon^{(j)}\bigr) \;-\; \frac{4j}{r^2}\Bigr| \;\le\; C\,\frac{\epsilon\,|\log\epsilon|}{r^2}.
\]
\end{corollary}

\begin{proof}
The min-max step of the preceding proof gives $\lambda_0^{\mathrm{ess}}(U_\epsilon^{(j)}) \ge (1 - C\epsilon|\log\epsilon|)\,4j/r^2$. Conversely, the transported quasi-modes $\Phi^{(j)}_n$ form, for $n \to \infty$ at fixed $\epsilon$, a Weyl-type sequence whose normalized residual $\|(-\Lap_{U_\epsilon^{(j)}} - 4j/r^2)\Phi^{(j)}_n\|_{L^2}/\|\Phi^{(j)}_n\|_{L^2}$ is bounded, in the limit, by $C'\epsilon|\log\epsilon|/r^2$; hence $\operatorname{dist}\bigl(4j/r^2,\,\sigma_{\mathrm{ess}}(-\Lap_{U_\epsilon^{(j)}})\bigr) \le C'\epsilon|\log\epsilon|/r^2$, and since the min-max bound excludes essential spectrum below $(1 - C\epsilon|\log\epsilon|)\,4j/r^2$, the essential-spectrum point so produced bounds $\lambda_0^{\mathrm{ess}}$ from above by $4j/r^2 + C'\epsilon|\log\epsilon|/r^2$. The two bounds combine into the display.
\end{proof}

\begin{remark}[On the strength of the metric-comparison rate]\label{rmk:rate}
The remainder rate in \eqref{eq:metric-product-bound} is $O(\epsilon|\log\epsilon|)$ rather than $O(\epsilon)$, owing to the logarithmic factor in the gauge-fixed position-position cross-block bound \eqref{eq:OD-horiz}; this is consistent with the logarithm appearing in the symmetric next-order remainder of Lemma~\ref{lem:cross-term} (cf.\ \cite{lemma-ct-higher-order}). The $O(\epsilon|\log\epsilon|) \to 0$ rate is amply sufficient for Theorem~\ref{thm:M_k-collar-conditional}. The strengthened $C^2$ form \eqref{eq:metric-product-bound-C2}, together with Lemma~\ref{app:lem:iterated-h}, is what makes the Mourre/Sahbani machinery of Theorem~\ref{thm:M_k-collar-spectral-type} applicable: an operator-norm (quadratic-form) bound alone would suffice for $\inf\sigma_{\mathrm{ess}}$ but would not control the iterated commutators required for the $C^{1,1}$-regularity hypothesis of \cite{Sahbani1997}.
\end{remark}

\begin{remark}\label{rmk:collar-vs-body}
The body of $\cM_k$ for $k \ge 2$ has curvature variation in the interior (away from collars), where neither McKean nor the asymptotic-product argument directly bounds the local spectrum. Theorem~\ref{thm:M_k-collar-conditional} is specifically a statement about the Uhlenbeck collar regions, where the geometry is asymptotically rigid.
\end{remark}

\subsection{Spectral type on the Uhlenbeck collar}\label{sec:mourre}

The metric-comparison estimate \eqref{eq:metric-product-bound} is sufficient to identify the bottom of the essential spectrum; with the iterated-commutator bounds of Appendix~\ref{app:weighted-sobolev}, \S\ref{app:sec:iterated}, it also identifies the \emph{spectral type} above the threshold via Mourre's commutator method. We use the formulation due to Sahbani \cite{Sahbani1997}, which improves the classical Amrein--Boutet de Monvel--Georgescu framework \cite{ABMG1996} to the optimal $C^{1,1}$-regularity class for the conjugate operator.

\begin{theorem}[Spectral type on Uhlenbeck collars]\label{thm:M_k-collar-spectral-type}
Let $U_\epsilon^{(j)} \subset \cM_k(S^4_r)$ be the codimension-$j$ Uhlenbeck collar as in Theorem~\ref{thm:M_k-collar-conditional}. For every $\delta > 0$, there exists $\epsilon_0 = \epsilon_0(\delta, j, r) > 0$ such that for all $0 < \epsilon < \epsilon_0$ the Laplace--Beltrami operator $-\Lap_{U_\epsilon^{(j)}}$ on the collar satisfies, on the spectral interval $I_\delta := (4j/r^2 + \delta, \infty)$:
\begin{enumerate}\setlength{\itemsep}{2pt}
\item[\textup{(i)}] A strict Mourre estimate holds with conjugate operator $A_j$ defined initially on $C_c^\infty(U_\epsilon^{(j)})$ by $A_j = \sum_{i=1}^j \tfrac12(y_i\partial_{y_i} + \partial_{y_i} y_i)\chi(y_i \ge Y_0)$ (with cutoff $\chi$ supported in the cusp ends), and then extended by closure to an essentially self-adjoint operator on $L^2(U_\epsilon^{(j)})$ (essential self-adjointness follows from the standard theory of first-order symmetric operators generating one-parameter unitary groups: Nelson's analytic vector theorem, Reed--Simon \cite[Thm.~X.39]{ReedSimonII}, with $C_c^\infty$ as the core of analytic vectors for the cusp dilation flow):
\[
E_{I_\delta}(-\Lap_{U_\epsilon^{(j)}})\,[-\Lap_{U_\epsilon^{(j)}},\,iA_j]\,E_{I_\delta}(-\Lap_{U_\epsilon^{(j)}})\;\ge\;\delta\,E_{I_\delta}(-\Lap_{U_\epsilon^{(j)}})\;+\;K_\epsilon,
\]
with $K_\epsilon$ a compact remainder.
\item[\textup{(ii)}] $A_j$ is of class $C^{1,1}(-\Lap_{U_\epsilon^{(j)}})$ in the sense of \cite[Def.~6.2.2]{ABMG1996}; equivalently, both the first and second iterated commutators $[-\Lap_{U_\epsilon^{(j)}}, iA_j]$, $[[-\Lap_{U_\epsilon^{(j)}}, iA_j], iA_j]$ extend from $C_c^\infty$ to bounded operators $\mathcal{D}(-\Lap_{U_\epsilon^{(j)}})\to \mathcal{D}(-\Lap_{U_\epsilon^{(j)}})^*$ with norm $O(1)$ as $\epsilon\to 0$.
\item[\textup{(iii)}] Consequently \cite[Thm.~0.1]{Sahbani1997}, $-\Lap_{U_\epsilon^{(j)}}$ has no singular continuous spectrum on $I_\delta$; its point spectrum on $I_\delta$ is locally finite (no accumulation interior to $I_\delta$); and a global limiting absorption principle holds on $I_\delta$, namely the boundary values $\langle A_j\rangle^{-1/2}(-\Lap_{U_\epsilon^{(j)}} - \lambda \mp i0)^{-1}\langle A_j\rangle^{-1/2}$ exist as norm-continuous functions of $\lambda \in I_\delta$.
\end{enumerate}
\end{theorem}

\begin{proof}
For (i): work first on the exact product cusp $\cM_{k-j}\times \bigtimes_{i=1}^j (\bbH^5_r\text{-cusp})$. We follow the Froese--Hislop framework \cite{FroeseHislop1989}: the Mourre identity does \emph{not} hold for the unconjugated upper-half-space Laplacian as an operator identity (direct calculation gives $[-\Lap_{\bbH^5_r},\,iA_i^{(0)}] = 2i r^{-2} y_i^2\,\partial_{x^{(i)}}^2$, a transverse-Laplacian-valued operator that vanishes only on the radial sector); the clean Mourre identity instead lives on the \emph{conjugated radial Schr\"odinger model} obtained by spherical-harmonic decomposition on the horosphere and the standard half-density conjugation.

Concretely, on the $i$-th cusp factor in upper-half-space coordinates $(x^{(i)},y_i)$ with $x^{(i)}\in\bbR^4$, $y_i>0$, decompose $L^2(\bbH^5_r\text{-cusp})$ via Fourier transform in the horospherical coordinate $x^{(i)}\in\bbR^4$ (a flat $\bbR^4$ at infinity, on which the transverse Laplacian $-\Lap_{x^{(i)}}$ has continuous spectrum $[0,\infty)$). Each Fourier mode at $|k|^2 = \mu \ge 0$ contributes a nonnegative additive shift $\mu\cdot y_i^2/r^2$ to the radial Schr\"odinger operator below; the $\mu = 0$ (zero-mode) sector is the one we focus on, where the radial reduction is clean. For $\mu > 0$ the additive shift is positive and short-range as $y_i \to \infty$ (proportional to $y_i^2 \mu$ in the conjugated coordinate $u_i = \log y_i$, becoming $e^{2u_i}\mu \to \infty$ at $u_i \to +\infty$), strictly raising the bottom of essential spectrum on those modes, so the threshold $4/r^2$ is set by the zero-mode sector. On the $\ell$-th sector the Laplacian, conjugated by the half-density factor $y_i^{(n-1)/2} = y_i^{2}$ (with $n = 5$) and the substitution $u_i = \log y_i$, becomes a half-line Schr\"odinger operator on $L^2(\bbR, du_i)$:
\begin{equation}\label{eq:Hconj-ell}
H_{\mathrm{conj}}^{(\ell)} \;=\; -\partial_{u_i}^2 \;+\; \frac{V_\ell(u_i)}{r^2}, \qquad V_0(u_i) \equiv 4, \qquad V_\ell(u_i) = 4 + \ell(\ell+3)\,e^{-2u_i} \quad (\ell\ge 1).
\end{equation}
The constant $4/r^2$ is the McKean threshold; the Casimir term $\ell(\ell+3)e^{-2u_i}/r^2$ is short-range as $u_i\to+\infty$ and provides repulsion on the cusp. On each conjugated sector the dilation generator $A_{u_i} := \tfrac12(u_i D_{u_i} + D_{u_i} u_i)$ with $D_{u_i} = -i\partial_{u_i}$ produces the \emph{exact} Mourre identity on the $\ell=0$ sector,
\begin{equation}\label{eq:Hconj-mourre}
[H_{\mathrm{conj}}^{(0)},\, iA_{u_i}] \;=\; 2(-\partial_{u_i}^2) \;=\; 2\bigl(H_{\mathrm{conj}}^{(0)} - 4/r^2\bigr),
\end{equation}
verified symbolically in \cite{mourre-cusp-script} (this is the standard half-line dilation--Laplacian commutator). For $\ell\ge 1$ the Casimir term breaks the clean identity but adds a strictly positive contribution: $[H_{\mathrm{conj}}^{(\ell)},\, iA_{u_i}] = 2(-\partial_{u_i}^2) - r^{-2}\ell(\ell+3)\cdot 2u_i e^{-2u_i}$, and on the spectral subspace above the threshold $4/r^2 + \delta$ the short-range Casimir contribution is a relatively compact perturbation of $H_{\mathrm{conj}}^{(\ell)}$ on the cusp end, absorbed into $K_\epsilon$. Transferring back through the unitary conjugation, this gives on the $i$-th \emph{full} cusp factor the operator identity (modulo a compact remainder from the half-density conjugation and the spherical-harmonic decomposition, both controlled by the compact transverse slice)
\begin{equation}\label{eq:cusp-mourre}
[-\Lap_{\bbH^5_r\text{-cusp},i},\,iA_i^{(0)}] \;=\; 2\bigl(-\Lap_{\bbH^5_r\text{-cusp},i} - 4/r^2\bigr) + K^{(i)}_{\mathrm{cpt}},\qquad A_i^{(0)} := \tfrac12(y_i\partial_{y_i}+\partial_{y_i}y_i),
\end{equation}
where $A_i^{(0)}$ is the pullback of $A_{u_i}$ under $u_i = \log y_i$ (the half-density-conjugated form), and $K^{(i)}_{\mathrm{cpt}}$ is the relatively compact transverse-and-Casimir contribution. Summing over $i$ and using the tensor-product decomposition $-\Lap_{\mathrm{prod}} = -\Lap_{\cM_{k-j}}\otimes 1 + \sum_i 1\otimes(-\Lap_{\text{cusp},i})$:
\begin{equation}\label{eq:prod-mourre}
[-\Lap_{\mathrm{prod}},\,iA_j^{(0)}] \;=\; 2\bigl(-\Lap_{\mathrm{prod}} - 4j/r^2\bigr) - 2(-\Lap_{\cM_{k-j}})\otimes 1\;\ge\;2\bigl(-\Lap_{\mathrm{prod}} - 4j/r^2\bigr) - 2(-\Lap_{\cM_{k-j}}).
\end{equation}
The cutoff $\chi(y_i\ge Y_0)$ introduces a localization error supported on the compact transverse slice $\{y_i = Y_0\}$, which is a relatively compact perturbation of $-\Lap_{\mathrm{prod}}$ (the inclusion $H^1(\text{slice})\hookrightarrow L^2(\text{slice})$ is compact by Rellich), absorbed into the compact remainder $K_\epsilon$.

\smallskip
\noindent\emph{Residual factor contribution.} The conjugate operator $A_j = \sum_i \tfrac12(y_i\partial_{y_i}+\partial_{y_i}y_i)\chi(y_i\ge Y_0)$ acts \emph{only} on the cusp factors and is the identity on the residual factor $\cM_{k-j}(S^4_r)$. In the tensor decomposition $-\Lap_{\mathrm{prod}} = -\Lap_{\cM_{k-j}}\otimes 1 + 1\otimes(-\Lap_{\mathrm{bubbles}})$, the residual Laplacian therefore commutes with $A_j$:
\[
[-\Lap_{\cM_{k-j}}\otimes 1,\, iA_j] \;=\; 0
\]
identically (not "modulo a compact remainder"). The full commutator $[-\Lap_{\mathrm{prod}}, iA_j] = [1\otimes(-\Lap_{\mathrm{bubbles}}), iA_j]$ therefore lives entirely on the cusp factors, where the Froese--Hislop identity \eqref{eq:cusp-mourre} gives $\sum_i 2(1\otimes(-\Lap_{\mathrm{cusp},i}-4/r^2)) + \sum_i K_{\mathrm{cpt}}^{(i)}\otimes 1$. Tensoring with the residual identity, the contribution on the joint spectral subspace $E_{I_\delta}(-\Lap_{\mathrm{prod}})$ is exactly $2(-\Lap_{\mathrm{prod}}|_{\mathrm{cusp}} - 4j/r^2)$ on the cusp factors, which on $E_{I_\delta}$ is $\ge 2\delta$ on the cusp-energy part, plus the cumulative compact remainder $K_\epsilon = \sum_i K_{\mathrm{cpt}}^{(i)}$.

The metric perturbation $g_{U_\epsilon^{(j)}} = g_{\mathrm{prod}} + h$ with $\|h\|_{g_{\mathrm{prod}}}\le C\epsilon|\log\epsilon|$ produces the comparison
\[
[-\Lap_{U_\epsilon^{(j)}},\,iA_j]\;=\;[-\Lap_{\mathrm{prod}},\,iA_j^{(0)}] \;+\; R_\epsilon^{(1)},
\]
with $\pm R_\epsilon^{(1)} \le C\epsilon|\log\epsilon|\,(-\Lap_{\mathrm{prod}}+1)$ as quadratic forms, by Lemma~\ref{app:lem:iterated} below applied at order $n=1$. For $\epsilon < \epsilon_0(\delta)$ small enough this lower-bounds the Mourre constant by $\delta$ on $I_\delta$ as claimed.

For (ii): the second-iterated commutator
\[
[[-\Lap_{U_\epsilon^{(j)}},\,iA_j],\,iA_j]\;=\;[[-\Lap_{\mathrm{prod}},\,iA_j^{(0)}],\,iA_j^{(0)}]\;+\;R_\epsilon^{(2)}
\]
is bounded relative to $-\Lap_{\mathrm{prod}}+1$ with norm $O(1)$: in the conjugated radial Schr\"odinger frame \eqref{eq:Hconj-ell}, the principal symbol of each cusp sector is $-\partial_{u_i}^2$ and the half-line dilation commutator iterates exactly to $[[-\partial_{u_i}^2, iA_{u_i}], iA_{u_i}] = 4(-\partial_{u_i}^2)$, so the unperturbed second commutator equals $4(-\Lap_{\mathrm{prod}}-4j/r^2)$ modulo a relatively compact remainder (short-range Casimir derivatives and the spherical-harmonic compactness, as in (i)); $R_\epsilon^{(2)} = O(\epsilon|\log\epsilon|^2)$ by Lemma~\ref{app:lem:iterated} at $n=2$. This is the $C^{1,1}(-\Lap_{U_\epsilon^{(j)}})$ regularity needed by Sahbani \cite[Hyp.~H1--H3, p.~411]{Sahbani1997}; in fact one obtains the strictly stronger $C^2(-\Lap_{U_\epsilon^{(j)}})$ regularity.

Statement (iii) is then the conclusion of \cite[Thm.~0.1]{Sahbani1997} (or equivalently \cite[Thm.~7.5.4]{ABMG1996} together with the standard absence of $\sigma_{\mathrm{sc}}$ corollary).
\end{proof}

\begin{remark}[On the absolute-continuity strengthening]\label{rmk:AC}
The conjunction of statements (i)--(iii) of Theorem~\ref{thm:M_k-collar-spectral-type} implies that the spectrum of $-\Lap_{U_\epsilon^{(j)}}$ on $I_\delta$ is the union of (a) absolutely continuous spectrum filling $I_\delta$ and (b) at most finitely many embedded eigenvalues on $I_\delta$ for each compact sub-interval (locally finite point spectrum). Whether embedded eigenvalues are present at all on $I_\delta$ is not settled by Mourre theory and depends on the residual moduli factor; on the model cusp $\bbH^5_r$ itself the spectrum is purely absolutely continuous and no eigenvalues appear.
\end{remark}

\begin{remark}[$\mathrm{SU}(N)$ extension]\label{rmk:SUN}
Theorem~\ref{thm:M_k-collar-conditional} holds verbatim for the moduli space $\cM_k(S^4_r; \mathrm{SU}(N))$ of charge-$k$ anti-self-dual $\mathrm{SU}(N)$ connections, $N \ge 2$, with the same collar essential bottom $4j/r^2$. The $\eta$-tensor identity used in Lemma~\ref{lem:cross-term} operates entirely inside the $\mathfrak{su}(2)$-subalgebra in which each bubble is framed, and the canonical normalization $\mathrm{Tr}(T_a T_b) = \tfrac12\delta_{ab}$ makes the SU(N) common-subgroup cross-term identical to \eqref{eq:lemma-asymptotic}. Bubbles framed in distinct SU(2) subgroups of SU(N), related by $U \in \mathrm{SU}(N)$, acquire a framing-overlap scalar
\[
o(U) \;:=\; \tfrac{2}{3}\sum_{a=1}^{3} \mathrm{Tr}_{\mathrm{fund}}\bigl(T_a\,U\,T_a\,U^{-1}\bigr) \;\in\; [-1/3,\,1],
\]
where $\{T_1, T_2, T_3\}$ is the fixed basis of the reference $\mathfrak{su}(2) \subset \mathfrak{su}(N)$ subalgebra in which the first bubble is framed (with the canonical normalization $\mathrm{Tr}_{\mathrm{fund}}(T_a T_b) = \tfrac12\delta_{ab}$), and $U$ is the relative framing rotation to the second bubble's $\mathfrak{su}(2)$ subalgebra. The factor $2/3$ normalizes $o$ so that the same-subgroup case $U = 1$ gives $o(1) = (2/3)\cdot 3 \cdot \tfrac12 = 1$, recovering the SU(2) value; the value $-1/3$ is attained on the symmetric self-conjugate fundamental representation. In all cases the factor $o(U)$ preserves the $O(s_1 s_2/R^2)$ decay rate. The additional compact framing factors $\mathrm{SU}(N)/S(U(N{-}2)\times U(2))$ on each bubble are positively curved and contribute only discrete spectrum starting at $0$ to each cusp; they do not lower the essential bottom. Symbolic verification of the SU(N) cross-term and numerical confirmation at $N=2,3,4$ are recorded in \cite{lemma-ct-SUN}.
\end{remark}

%-- moved to companion supplementary note: SO(5)-isotypic reduction (Proposition, conditional on V_true matching the McKean cusp asymptote) --
\iffalse
\subsection{$\mathrm{SO}(5)$-isotypic reduction of the global bottom on $\cM_2(S^4_r)$}\label{sec:isotypic}

The collar bound of Theorem~\ref{thm:M_k-collar-conditional} controls the essential-spectrum contribution from each Uhlenbeck stratum but does not determine the global $\lambda_0(\cM_k(S^4_r))$ for $k \ge 2$. For $k=2$ we reduce the global question to a single $\mathrm{SO}(5)$-isotypic component: the trivial (invariant) component. On every non-trivial $\mathrm{SO}(5)$-isotypic component the Casimir potential generated by the principal-orbit volume forces the bottom to coincide with the McKean essential threshold $4/r^2$, with no $L^2$ eigenvalues below.

\begin{theorem}[$\mathrm{SO}(5)$-isotypic reduction for $\cM_2(S^4_r)$]\label{thm:SO5-isotypic}
Let $H = -\Lap$ act on $L^2(\cM_2(S^4_r))$ with the natural $L^2$ metric, and let $H = \bigoplus_{\ell\ge 0} H_\ell$ be the orthogonal decomposition into $\mathrm{SO}(5)$-isotypic components, with $\ell$ indexing the highest weights of the $\mathrm{SO}(5)$-representations and $H_0$ the $\mathrm{SO}(5)$-invariant component. Then for every $\ell \ge 1$,
\[
\inf\sigma(H_\ell) \;=\; \frac{4}{r^2},
\]
realized as the bottom of continuous spectrum, with $\sigma_{\mathrm{pp}}(H_\ell)\cap (0, 4/r^2) = \emptyset$. Consequently $\lambda_0(\cM_2(S^4_r)) = \min\{4/r^2,\ \inf\sigma(H_0)\}$, and the global bottom is realized in the $\mathrm{SO}(5)$-invariant subspace.
\end{theorem}

\begin{proof}[Proof sketch]
On the cohomogeneity-one slice associated to the $\mathrm{SO}(5)$-action, the radial Schr\"odinger operator on each $(\ell, 0)$-isotypic component (highest-weight $\ell$ symmetric representation of $\mathrm{SO}(5)$) is
\[
H_\ell \;=\; -\partial_s^2 \;+\; V_{\mathrm{true}}(s) \;+\; \frac{C_\ell}{r_{\mathrm{orbit}}(s)^2},\qquad C_\ell = \ell(\ell+3),
\]
acting on $L^2((0,\infty), J(s)\,ds)$ where $J(s)$ is the codimension-1 volume density and $V_{\mathrm{true}}(s) = \tfrac14(J'/J)^2 + \tfrac12(J'/J)'$ is the Liouville potential. The principal-orbit radius interpolates between $r_{\mathrm{orbit}}(s) \sim s$ at the orbifold tip $s\to 0$ and $r_{\mathrm{orbit}}(s) \sim r\sinh(s/r)$ at the cusp $s\to\infty$ \cite{m2-isotypic-decomposition}. At the cusp end, $V_{\mathrm{true}}(s) \to 4/r^2$ (McKean rate on the $\bbH^5_r$ horocycle, independently of the bare-vs-projected normalization discussion of \cite{m2-closure-higher-order}) and $C_\ell/r_{\mathrm{orbit}}(s)^2 \to 0$ exponentially; at the tip, the Casimir potential dominates as $C_\ell/s^2$, repulsive. By the Bargmann criterion (Reed--Simon IV \cite{ReedSimonIV}, Thm.~XIII.9),
\[
N_-(H_\ell - 4/r^2) \;\le\; \int_0^\infty s\,(V_{\mathrm{true}}(s) + C_\ell/r_{\mathrm{orbit}}(s)^2 - 4/r^2)_-\,ds.
\]
The Casimir contribution is strictly nonnegative everywhere and dominates at small $s$, so the negative part of the integrand vanishes on $(0, s_0]$ for some $s_0 > 0$ depending only on $\ell$. On $(s_0,\infty)$, $C_\ell/r_{\mathrm{orbit}}^2$ decays exponentially while $V_{\mathrm{true}}-4/r^2$ approaches $0$ as $s\to\infty$ by the McKean asymptote. Numerical evaluation \cite{m2-isotypic-decomposition} of the Bargmann integrand under the McKean cusp asymptote gives
\[
B_\ell \;:=\; \int_0^\infty s\,(V_{\mathrm{true}}+ C_\ell/r_{\mathrm{orbit}}^2 - 4/r^2)_-\,ds \;\approx\; 0.34,\ 0.26,\ 0.21,\ldots
\]
for $\ell = 1,2,3,\ldots$, monotonically decreasing in $\ell$, all strictly less than $1$. Hence $N_-(H_\ell - 4/r^2) = 0$ for every $\ell \ge 1$. The bottom $4/r^2$ is then the essential threshold of $H_\ell$, attained in continuous spectrum without bound states. The robustness of the Casimir-repulsion mechanism (uniform across admissible $V_{\mathrm{true}}$ matching the prescribed endpoint asymptotes, in contrast with the family-dependence at $\ell = 0$) is documented in \cite{m2-isotypic-decomposition}.
\end{proof}

\begin{remark}[Reduction of the global question to $\ell = 0$]\label{rmk:isotypic-reduction}
Theorem~\ref{thm:SO5-isotypic} reduces determination of $\lambda_0(\cM_2(S^4_r))$ to the single problem of computing $\inf\sigma(H_0)$ on the $\mathrm{SO}(5)$-invariant subspace. On $\ell = 0$ the Casimir-repulsion mechanism does not apply (since $C_0 = 0$) and the question reduces to a one-dimensional Bargmann problem on the radial Schr\"odinger operator with the true Liouville potential $V_{\mathrm{true}}(s)$, which is not analytically closed by the methods of this note; see the discussion in the companion supplementary note.
\end{remark}

\subsection{Finiteness of point spectrum below threshold and absence of singular continuous spectrum}\label{sec:MM-finiteness}

A separate and complementary spectral statement uses the asymptotic-hyperbolic structure of the Uhlenbeck collars (Lemma~\ref{lem:OD}, Theorem~\ref{thm:M_k-collar-conditional}) together with the Mazzeo--Melrose 0-calculus on asymptotically hyperbolic manifolds. The collars satisfy the Mazzeo--Melrose hypotheses up to a polyhomogeneous correction with index set $\subseteq \{(j,l) : j \ge 1,\ 0 \le l \le 1\}$, owing to the $O(\epsilon|\log\epsilon|)$ rate in \eqref{eq:metric-product-bound}, which is within the allowable class of the modern 0-calculus extension (Mazzeo \cite{Mazzeo1991}, Vasy \cite{Vasy2013}). The stratified structure of the conformal boundary (an iterated Uhlenbeck stratification of $\partial\cM_k$) is handled by the stratified / iterated-edge extension (Albin--Leichtnam--Mazzeo--Piazza \cite{ALMP2012}).

\begin{theorem}[Finiteness of point spectrum and absence of $\sigma_{\mathrm{sc}}$]\label{thm:MM-finiteness}
Let $\cM_k(S^4_r)$, $k \ge 1$, be equipped with the natural $L^2$ metric, and let $-\Lap$ be its Friedrichs $L^2$-Laplacian. Then
\begin{enumerate}
\item[(a)] $\sigma_{\mathrm{ess}}(-\Lap) = [4/r^2, \infty)$;
\item[(b)] $\sigma_{\mathrm{pp}}(-\Lap)\cap (0, 4/r^2)$ is a finite set (possibly empty), each eigenvalue of finite multiplicity;
\item[(c)] $\sigma_{\mathrm{sc}}(-\Lap)\cap(0, \infty) = \emptyset$;
\item[(d)] the resolvent $(-\Lap - z)^{-1}: C_c^\infty \to C^\infty$ admits a meromorphic continuation from $\{\Im z > 0\}$ across the cut $[4/r^2,\infty)$ into a logarithmic cover, with poles in the physical sheet corresponding exactly to the $L^2$ eigenvalues of (b).
\end{enumerate}
\end{theorem}

\begin{proof}[Proof sketch]
Statement (a) follows from Theorem~\ref{thm:M_k-collar-conditional} (upper bound) together with the Persson-type/Donnelly--Li tail-comparison lemma applied to the asymptotic-product cusps (lower bound; cf.\ Donnelly \cite{Donnelly1981}). Statements (b), (c), (d) follow from the 0-calculus parametrix construction for the resolvent on asymptotically hyperbolic manifolds (Mazzeo--Melrose \cite{MazzeoMelrose1987}, Mazzeo \cite{Mazzeo1991} Thm.~7.1; cf.\ Guillarmou \cite{Guillarmou2005}, Vasy \cite{Vasy2013}): the resolvent off the threshold is a $0$-pseudodifferential operator of order $-2$ with polyhomogeneous Schwartz kernel at the front face, modulo compact errors. The $O(\epsilon|\log\epsilon|)$ metric remainder in \eqref{eq:metric-product-bound} lies in the allowable polyhomogeneous index family (index set $\subset \{(j,l) : j \ge 1, 0 \le l \le 1\}$), and the stratified Uhlenbeck-boundary structure is handled by the iterated-edge calculus of Albin--Leichtnam--Mazzeo--Piazza \cite{ALMP2012}. Compactness of the embedding $\mathrm{Dom}(\Lap^{1/2}) \hookrightarrow L^2_{\mathrm{loc}}$ relative to the essential threshold gives the finiteness of $\sigma_{\mathrm{pp}}\cap(0, 4/r^2)$; absence of $\sigma_{\mathrm{sc}}$ follows from the meromorphic resolvent continuation combined with Mourre theory (Theorem~\ref{thm:M_k-collar-spectral-type}) on the cusp ends. The two routes are consistent: Mourre theory rules out $\sigma_{\mathrm{sc}}$ above threshold; the 0-calculus rules it out everywhere.
\end{proof}

\begin{remark}[Caveats on the citation chain]\label{rmk:MM-caveats}
Statement (a) is fully proved by Theorem~\ref{thm:M_k-collar-conditional} + Persson. Statements (c) above threshold and the absence-of-embedded-eigenvalues part of (c) follow from Theorem~\ref{thm:M_k-collar-spectral-type}. Statements (b) and (d), together with $\sigma_{\mathrm{sc}}\cap(0, 4/r^2) = \emptyset$, are extracted via the 0-calculus citation chain just sketched: the load-bearing technical input is the iterated-edge $0$-calculus parametrix construction for stratified asymptotically hyperbolic manifolds, which is by now standard but is not a one-line citation. A self-contained derivation specialized to $\cM_k(S^4_r)$ remains a worthwhile project. The finiteness statement (b) does \emph{not} on its own settle whether $\sigma_{\mathrm{pp}}\cap (0, 4/r^2) = \emptyset$ (i.e., whether $\lambda_0(\cM_k(S^4_r)) = 4/r^2$): see the discussion in the companion supplementary note for the case $k=2$ and the open question for $k\ge 3$.
\end{remark}
\fi
%-- end moved content --

\section{Remarks and open questions}

\begin{remark}[Sign of $\mathrm{Ric}$ on instanton moduli]
$\bbH^5_r$ has $\mathrm{Ric} = -4g/r^2$, strictly negative. The Bakry--\'Emery--Lichnerowicz framework with positive Ricci would give a discrete eigenvalue gap; with negative Ricci, it gives a continuous spectrum with positive bottom, consistent with Theorem~\ref{thm:M1-spectrum} and Theorem~\ref{thm:M_k-collar-conditional}. Singer's orbit-space proposal \cite{Singer1981} requires positive Ricci on the relevant sector for Lichnerowicz--Obata to give a spectral gap; the instanton sectors carry continuous spectrum bounded by McKean rather than discrete eigenvalues.
\end{remark}

\begin{remark}[Open: global $\lambda_0(\cM_k)$ for $k \ge 2$]
Theorem~\ref{thm:M_k-collar-conditional} bounds the essential-spectrum bottom from each collar but does not determine the global $\lambda_0(\cM_k)$ for $k \ge 2$. The natural attack uses the explicit ADHM moduli geometry and the Groisser--Parker interior curvature bounds, combined with a Bargmann-type argument on the radial reduction in the SO(5)-invariant cohomogeneity-one sector.
\end{remark}

\begin{remark}[Additional structural results in the companion supplementary note]\label{rmk:companion}
Two structural reductions of the global $\lambda_0(\cM_2)$ question, an $\mathrm{SO}(5)$-isotypic decomposition reducing the question to the invariant subspace via Bargmann + Casimir repulsion on the $\ell \ge 1$ components, and a Mazzeo--Melrose-type finiteness-and-meromorphic-resolvent statement for $\sigma_{\mathrm{pp}} \cap (0, 4/r^2)$ via iterated-edge $0$-calculus on stratified asymptotically hyperbolic manifolds, are recorded as conditional propositions in the companion supplementary note. Each rests on inputs (uniform-in-family Bargmann robustness for the $\mathrm{SO}(5)$-isotypic reduction; stratified iterated-edge calculus verification for the $0$-calculus statements) whose self-contained derivations are beyond the scope of this note.
\end{remark}

\appendix

\section{Weighted-Sobolev proof of Lemma~\ref{lem:OD}(iv)--(v)}\label{app:weighted-sobolev}

This appendix proves the two off-diagonal cross-block bounds of Lemma~\ref{lem:OD} that were stated in the body without a self-contained calculation. The argument is the standard Taubes \cite{Taubes1988}/Donaldson--Kronheimer \cite{DK1990}~\S 7.3 weighted-Sobolev framework, with the scale-uniformity made explicit. We cite Bartnik \cite{Bartnik1986} and Lockhart--McOwen \cite{LMcO1985} for the abstract Fredholm theory of the weighted Laplacian on $\bbR^n$; everything else is proved here.

\subsection{Weighted Sobolev setup}\label{app:sec:setup}

On $\bbR^4$, fix a smooth weight $\sigma(x) = (1+|x-x_0|^2)^{1/2}$ and for $\beta \in \bbR$, define
\[
\|f\|_{L^2_\beta}^2 \;:=\; \int_{\bbR^4} |f(x)|^2\,\sigma(x)^{-2\beta - 4}\,d^4x,
\qquad
\|f\|_{W^{2,2}_\beta}^2 \;:=\; \sum_{|\alpha|\le 2}\|\sigma^{|\alpha|}\nabla^\alpha f\|_{L^2_\beta}^2 ,
\]
following the Bartnik convention \cite{Bartnik1986}: a function with pointwise rate $|f(x)| \le C\,\sigma(x)^{\beta'}$ lies in $L^2_\beta$ for $\beta' < \beta$, and the borderline rate $|f|\sim \sigma^{\beta'}$ is in $L^2_\beta$ iff $\beta' < \beta$ (strict). The pairing $L^2_\beta \times L^2_{-4-\beta} \to \bbR$ identifies $(L^2_\beta)^* \cong L^2_{-4-\beta}$ \cite{Bartnik1986}. \emph{Convention note:} this is the Bartnik sign convention; the Lockhart--McOwen weight convention \cite{LMcO1985} uses the opposite sign for $\beta$. Both yield the same Fredholm theory after a sign flip; we work throughout with the Bartnik convention, in which $\Lap: W^{2,2}_\beta \to L^2_{\beta - 2}$ on $\bbR^4$ is Fredholm for non-critical $\beta$ (the critical exponents are the integers $\beta \in \{-2, -3, -4, \dots\}$ corresponding to the kernel/cokernel obstructions of the Laplacian on the asymptotically Euclidean half-line).

We work with the weight $\beta = -3/2$. This is the natural weight for the BPST tangent space because:
\begin{itemize}\setlength{\itemsep}{2pt}
\item The field-strength rate $|F_A(x)| = O(s^2/|x-x_0|^4)$ lies in $L^2_\beta$ for any $\beta > -4$.
\item The bare position-derivative tail $|v_{x,\rho}(x)| = O(|x-x_0|^{-2})$ fails to lie in plain $L^2 = L^2_{-2}$ (it sits exactly at the critical exponent $\beta_\mathrm{crit} = -2$); however, by the Bartnik strict-inequality criterion above, $v_{x,\rho} \in L^2_\beta$ for every $\beta > -2$, in particular for the working weight $\beta = -3/2$, with norm growing like $|\log s|^{1/2}$ as $s \to 0$. The Coulomb obstruction is therefore that $v_{x,\rho}$ does not lie in the \emph{energy space} (in which derivatives are also weighted) without renormalization. The horizontal projection $\hat v_{x,\rho}$ instead satisfies $|\hat v_{x,\rho}| = O(s^2/(r+s)^3)$ (Lemma~\ref{app:lem:tail}) and belongs to $L^2_{-3/2}$ with norm uniformly bounded in $s$.
\item The horizontal-projected rate $|\hat v_{x,\rho}(x)| = O(s^2/|x-x_0|^3)$ is strictly inside $L^2_{-3/2}$, with $L^2_\beta$-norm uniformly bounded in $s$.
\end{itemize}

On the BPST background $A$ with scale $s$, write $\Lap_A = d_A^* d_A$ acting on $\fg$-valued functions. The conformal/scale invariance of the BPST connection under $x \mapsto x_0 + s\, x'$, $A \mapsto s^{-1} A'$ rescales $\Lap_A \mapsto s^{-2}\Lap_{A'}$, and rescales the weighted norms by $\|f\|_{L^2_\beta} \mapsto s^{2+\beta}\|f'\|_{L^2_\beta}$ (the factor $s^{2}$ from the volume Jacobian, plus $s^\beta$ from the weight). The combination is what produces the scale-uniform constants below.

\begin{lemma}[Scale-uniform weighted Coulomb estimate]\label{app:lem:coulomb}
There exists $C_0$, independent of the BPST scale $s$, such that for every $\fg$-valued $1$-form $v$ on $\bbR^4$ with $d_A^* v \in L^2_{-1/2}$, the equation
\[
\Lap_A \phi \;=\; d_A^* v
\]
admits a unique solution $\phi \in L^2_{-3/2}$ satisfying
\[
\|\phi\|_{W^{2,2}_{-3/2}} \;\le\; C_0\,\|d_A^* v\|_{L^2_{-1/2}}.
\]
\end{lemma}

\begin{proof}
For $A \equiv 0$ (the flat background), the weight $\beta = -3/2$ is non-critical for $\Lap$ on $\bbR^4$ in the sense of Lockhart--McOwen \cite{LMcO1985}, lying strictly between the critical exponents $\beta = -2$ (kernel: constants) and $\beta = -4$ (cokernel obstruction). Bartnik's theorem \cite[Thm.~1.7]{Bartnik1986} then states that $\Lap: W^{2,2}_{-3/2} \to L^2_{-1/2}$ is an isomorphism, giving an estimate $\|\phi\|_{W^{2,2}_{-3/2}}\le C_1\|\Lap\phi\|_{L^2_{-1/2}}$ with $C_1$ an absolute constant. \emph{Convention reconciliation:} the Bartnik convention adopted here pairs a domain weight $\beta = -3/2$ with target weight $-1/2$ via the duality $(L^2_{-1/2})^* \cong L^2_{-7/2}$; the equivalent alternative convention $\Lap: W^{2,2}_\beta \to L^2_{\beta-2}$ yields target weight $-7/2$ at $\beta = -3/2$. Both conventions give identical Fredholm content; we work throughout with the Bartnik form.

For the BPST background at scale $s = 1$ centered at $x_0 = 0$, write $\Lap_A = \Lap + [A, \cdot]$ where the second term is a first-order operator with coefficient $|A(x)| = O(|x|/(1+|x|^2)) = O(\sigma^{-1})$ pointwise. The perturbation $K := \Lap_A - \Lap = 2A\cdot\nabla + (\nabla\cdot A) + A\cdot A$ has $K: W^{2,2}_{-3/2}\to L^2_{-1/2}$ compact (the coefficients decay like $\sigma^{-1}$, $\sigma^{-2}$, $\sigma^{-2}$ respectively, gaining one full unit of decay relative to the borderline). Triviality of the kernel of $\Lap_A$ on $\fg$-valued functions in $L^2_{-3/2}$ is the AHS/Taubes infinitesimal-deformation vanishing: any $\phi \in L^2_{-3/2}$ with $\Lap_A\phi = 0$ satisfies $\langle \phi, \Lap_A\phi\rangle_{L^2} = \|d_A\phi\|^2_{L^2} = 0$ (the integration by parts is justified by the strict $\sigma^{-3/2}$ decay), so $d_A\phi = 0$; on an irreducible bundle the only covariantly constant $\fg$-valued function is $\phi \equiv 0$ \cite[Prop.~3.4]{AHS1978} (see also \cite[Lem.~4.2.4]{DK1990}). Hence $\Lap_A: W^{2,2}_{-3/2}\to L^2_{-1/2}$ is invertible, with an estimate $\|\phi\|_{W^{2,2}_{-3/2}}\le C_0 \|\Lap_A \phi\|_{L^2_{-1/2}}$.

We now extract the $s$-uniformity by conformal scaling. Under the dilation $\Phi_s: x' \mapsto x_0+sx'$, the BPST connection at scale $s$ is the pullback (up to a constant gauge transformation) of the unit-scale BPST connection: $\Phi_s^*A_s = A_{s=1}$, so $\Lap_{A_s} = s^{-2}\,\Phi_s^*\Lap_{A_{s=1}}$. Track the scalings on each side of $\|\phi\|_{W^{2,2}_\beta(\bbR^4)} \le C\|d_A^*v\|_{L^2_{\beta+1}(\bbR^4)}$:
\begin{itemize}\setlength{\itemsep}{2pt}
\item Volume element: $d^4 x = s^4\, d^4 x'$.
\item Weight $\sigma(x) = (1 + |x|^2)^{1/2}$ at large $|x|$: $\sigma(x) \sim s\,\sigma(x')$, so $\sigma^{-2\beta-4}\,d^4x \sim s^{-2\beta-4}\cdot s^4\, d^4 x' = s^{-2\beta}\,d^4 x'$.
\item $L^2_\beta$ norm of $\phi$ scales by $s^{-\beta+0}$ (after taking square root): $\|\phi\|_{L^2_\beta(\bbR^4_x)} = s^{-\beta}\,\|\phi'\|_{L^2_\beta(\bbR^4_{x'})}$ where $\phi'(x') = \phi(\Phi_s(x'))$.
\item Two derivatives ($\partial^2$ adds factor $s^{-2}$ in the operator, while integration by parts in the norm-squared form gives a factor $s^{+2}$ extra in the norm): $\|\phi\|_{W^{2,2}_\beta} \asymp s^{-\beta - 2 + \text{volume}}\cdot \|\phi'\|_{W^{2,2}_\beta}$. The full bookkeeping gives $\|\phi\|_{W^{2,2}_\beta(\bbR^4_x)} = s^{-\beta+2}\,\|\phi'\|_{W^{2,2}_\beta(\bbR^4_{x'})}$ at $\beta = -3/2$.
\item Similarly $\|d_A^*v\|_{L^2_{\beta+1}}$ scales by $s^{-(\beta+1)+2} = s^{-\beta+1}$ (one derivative loses $s^{-1}$ relative to the norm; the additional $-1$ in the weight shift to $\beta+1$ adds another).
\end{itemize}
The two sides scale by the same total power $s^{-\beta+2} = s^{7/2}$ at $\beta = -3/2$, so the inequality $\|\phi\|_{W^{2,2}_\beta} \le C \|d_A^*v\|_{L^2_{\beta+1}}$ is invariant under $\Phi_s$, and the constant $C_0$ obtained at $s = 1$ transfers verbatim to all $s > 0$.
\end{proof}

\subsection{Tail decay of the horizontal projection}\label{app:sec:tail}

Let $A$ be a BPST instanton centered at $x_0$ with scale $s$, and let $v_{x,\rho} = \partial A/\partial x_0^\rho$ be the bare position-derivative tangent vector. We work in the regular gauge $A_\mu^a = 2\eta^a_{\mu\nu}(x-x_0)^\nu/D$, $D := |x-x_0|^2 + s^2$.

\begin{lemma}[Coulomb source and pointwise tail]\label{app:lem:tail}
The horizontal correction $\phi^\rho$ defined by $\Lap_A\phi^\rho = d_A^* v_{x,\rho}$ satisfies, with $r := |x-x_0|$,
\[
|\phi^\rho(x)| \le C\,\frac{s^2}{r+s},\qquad |d_A\phi^\rho(x)| \le C\,\frac{s^2}{(r+s)^2},
\]
and hence $\hat v_{x,\rho} := v_{x,\rho} - d_A\phi^\rho$ satisfies the field-strength rate
\[
|\hat v_{x,\rho}(x)| \le C\,\frac{s^2}{(r+s)^3}.
\]
In particular $\|\hat v_{x,\rho}\|_{L^2(\bbR^4)} \le C$ uniformly in $s$, and $\|\hat v_{x,\rho}\|_{L^2_\beta} \le C$ for any $\beta > -3/2$ (Bartnik weight).
\end{lemma}

\begin{proof}
A direct calculation in the regular gauge gives $d_A^* v_{x,\rho} = -2\eta^a_{\mu\rho}\partial^\mu(1/D)$ (the BPST identity $\Lap A = 0$ in this gauge implies the only surviving term comes from differentiating the explicit $x_0$-dependence in $D$). The pointwise rate is $|d_A^* v_{x,\rho}(x)| = O(s^2(r+s)/D^2) = O(s^2/(r+s)^3)$. Computing the weighted norm,
\[
\|d_A^* v_{x,\rho}\|_{L^2_{-1/2}}^2 = \int \frac{s^4}{(r+s)^6}\,(1+r^2)^{-3/2-2}\,d^4x \;=\; O(s^4)\!\int_0^\infty \frac{r^3\,dr}{(r+s)^6 (1+r)^{7}} \,\le\, C\,s^4\cdot s^{-4} = C
\]
(both inner ($r\sim s$) and outer regions contribute $O(1)$ after the $s^4$ prefactor; explicit scaling $r = s\rho$ inner, $r$ outer). Thus $d_A^*v_{x,\rho}\in L^2_{-1/2}$ uniformly in $s$, with norm $\|d_A^*v_{x,\rho}\|_{L^2_{-1/2}} = O(1)$.

Lemma~\ref{app:lem:coulomb} gives $\phi^\rho \in W^{2,2}_{-3/2}$ with $\|\phi^\rho\|_{W^{2,2}_{-3/2}}\le C_0\cdot O(1) = O(1)$, uniformly in $s$. The pointwise bound on $\phi^\rho$ now follows from the standard interior-elliptic-regularity bootstrap combined with the explicit form of the source: $d_A^*v_{x,\rho}$ is $O(s^2(r+s)^{-3})$, and convolution with the $\bbR^4$ Newton kernel $\sim |y|^{-2}$ produces $|\phi^\rho| = O(s^2/(r+s))$. Differentiating once gives $|d_A\phi^\rho| = O(s^2/(r+s)^2)$.

We now identify the cancellation explicitly. Translations $x_0 \mapsto x_0 + \epsilon e_\rho$ are gauge equivalences of the BPST family: the connection $A(x; x_0 + \epsilon e_\rho)$ is gauge-equivalent to $A(x; x_0)$ via a smooth gauge transformation $g_\epsilon \to 1$ as $\epsilon \to 0$, so the infinitesimal translation $v_{x,\rho} = \partial A/\partial x_0^\rho$ equals an infinitesimal gauge transformation $-d_A \psi^\rho$ \emph{up to a non-pure-gauge piece}. Writing the gauge transformation as $\psi^\rho \in \Omega^0(\mathbb{R}^4; \fg)$, the Bianchi identity $d_A F_A = 0$ implies that the non-pure-gauge piece $\hat v_{x,\rho} = v_{x,\rho} + d_A\psi^\rho$ contracts only with the field-strength tensor, i.e.\ inherits the curvature decay rate $|F_A| = O(s^2(r+s)^{-4})$ multiplied by a length factor. The Coulomb-projected $\phi^\rho$ is precisely the unique gauge transformation $\psi^\rho$ that makes $\hat v_{x,\rho}$ Coulomb-horizontal ($d_A^*\hat v_{x,\rho} = 0$); existence and uniqueness are Lemma~\ref{app:lem:coulomb}. Hence
\[
|\hat v_{x,\rho}(x)| = |v_{x,\rho} - d_A\phi^\rho| = O\bigl(s^2(r+s)^{-3}\bigr)
\]
pointwise, the leading $|x-x_0|^{-2}$ gauge tails of $v_{x,\rho}$ and of $d_A\phi^\rho$ canceling exactly. The $L^2$ norm is $\|\hat v_{x,\rho}\|_{L^2}^2 = O(s^4)\int r^3 dr/(r+s)^6 = O(1)$.
\end{proof}

\subsection{Green's-kernel pointwise estimate for the BPST horizontal projection}\label{app:sec:green}

The pointwise bound $|\hat v_{x,\rho}(x)| = O(s^2/(r+s)^3)$ of Lemma~\ref{app:lem:tail} was deduced from the Bianchi-identity / gauge-equivalence argument, which is structural. Following the requests of the external review, we record here the same bound via the explicit Green's-kernel representation of the BPST-background Laplacian, which makes the cancellation of the long-range $|x-x_0|^{-2}$ tail manifest. The argument is standard (Taubes \cite{Taubes1988} \S 4; Donaldson--Kronheimer \cite{DK1990} \S 7.3) but we spell it out for transparency.

\begin{lemma}[Green's-kernel pointwise estimate for the BPST horizontal projection]\label{app:lem:green}
Let $A$ be a BPST instanton on $\bbR^4$ centered at $x_0$ with scale $s$, and let $G_A(x,x')$ denote the integral kernel of the Green's operator $\Lap_A^{-1}$ for the BPST-background Laplacian on $\fg$-valued $0$-forms, in the Bartnik weight class $\beta = -3/2$ (Lemma~\ref{app:lem:coulomb}). Write $G_0(x,x') = (4\pi^2)^{-1}|x-x'|^{-2}$ for the bare $\bbR^4$ Newton kernel. Then:
\begin{enumerate}\setlength{\itemsep}{2pt}
\item[\textup{(a)}] $G_A(x,x') = G_0(x,x') + R_A(x,x')$ with the remainder satisfying the pointwise bound
\[
|R_A(x, x')| \;\le\; \frac{C}{(1+|x-x'|)^{2}\,(1+|x-x_0|)\,(1+|x'-x_0|)},
\]
uniformly in the scale $s$ in the Bartnik--Lockhart--McOwen weight class.
\item[\textup{(b)}] The horizontal correction $\phi^\rho = G_A\,d_A^*v_{x,\rho}$ satisfies the explicit kernel representation
\[
\phi^\rho(x) \;=\; \int_{\bbR^4} G_A(x, x')\,d_A^* v_{x,\rho}(x')\,d^4x',
\]
and the pointwise bound $|d_A\phi^\rho(x)| \le C\,s^2/(r+s)^2$ of Lemma~\ref{app:lem:tail} is obtained by direct integration against the explicit source $|d_A^*v_{x,\rho}(x')| = O(s^2/(|x'-x_0|+s)^3)$.
\item[\textup{(c)}] The leading $|x-x_0|^{-2}$ tail of the bare $v_{x,\rho}$ and the leading $|x-x_0|^{-2}$ tail of $d_A\phi^\rho$ cancel pointwise, leaving the field-strength rate
\[
|\hat v_{x,\rho}(x)| \;\le\; \frac{C\,s^2}{(|x-x_0|+s)^3}
\]
uniformly in $s$ and $x$. The constant $C$ depends only on the BPST profile.
\end{enumerate}
\end{lemma}

\begin{proof}
\emph{(a)} On $\bbR^4$ the bare Laplacian Green's function is $G_0(x,x') = (4\pi^2)^{-1}|x-x'|^{-2}$. The BPST background contributes a perturbation
\[
\Lap_A - \Lap_0 \;=\; -[A^\mu, \partial_\mu \cdot] - \partial^\mu[A_\mu, \cdot] - [A^\mu, [A_\mu, \cdot]],
\]
with coefficients of pointwise size $|A| = O(s/(r+s)^{2})$ and $|A|^2 = O(s^2/(r+s)^4)$. In the Bartnik weight $\beta = -3/2$ class, Lemma~\ref{app:lem:coulomb} establishes that $\Lap_A: W^{2,2}_{-3/2}\to L^2_{-1/2}$ is an isomorphism uniformly in $s$, so its inverse $\Lap_A^{-1}$ has an integral kernel $G_A$ admitting the resolvent series
\[
G_A \;=\; G_0 \;-\; G_0\,(\Lap_A - \Lap_0)\,G_0 \;+\; G_0\,(\Lap_A - \Lap_0)\,G_0\,(\Lap_A - \Lap_0)\,G_0 \;-\; \cdots,
\]
convergent in the weighted operator-norm topology by the scale-uniform Coulomb estimate. The first-order remainder is bounded by
\[
|R_A^{(1)}(x,x')| \;\le\; \int |G_0(x,y)|\,|A(y)|\,|\partial G_0(y,x')|\,d^4y \;+\; (\text{lower order}),
\]
and the weighted Riesz convolution estimate $\int |x-y|^{-2}(1+|y|)^{-2}(1+|y|)^{-2}|y-x'|^{-3}\,d^4y \le C(1+|x-x'|)^{-2}(1+|x|)^{-1}$ on $\bbR^4$, a Hardy--Littlewood--Sobolev / Riesz-composition statement (sum of homogeneity orders $2+3 = 5 > 4 = \dim$ giving the extra factor $(1+|x|)^{-1}$); see Stein \cite{Stein1970}, Ch.~V, Thm.~1, gives the pointwise bound. Higher-order terms gain further factors of $(1+|x-x_0|)^{-1}$, so they are subleading. The translation-invariance argument may be reformulated as: $G_A - G_0 = G_0\,V\,G_A$ with $V = \Lap_A - \Lap_0$, and a direct Schur test in the Bartnik weight closes the iteration.

\emph{(b)} Direct substitution: $\Lap_A \phi^\rho = d_A^* v_{x,\rho}$ inverts to $\phi^\rho(x) = \int G_A(x,x')\,d_A^*v_{x,\rho}(x')\,d^4x'$. Split $G_A = G_0 + R_A$:
\[
\phi^\rho(x) \;=\; \underbrace{\int G_0(x,x')\,d_A^*v_{x,\rho}(x')\,d^4x'}_{=: \phi_0^\rho(x)} \;+\; \int R_A(x,x')\,d_A^*v_{x,\rho}(x')\,d^4x'.
\]
The first integral is the classical $\bbR^4$ Newton-kernel convolution of the source $d_A^* v_{x,\rho}$, which by the regular-gauge identity $d_A^* v_{x,\rho} = -2\eta^a_{\mu\rho}\partial^\mu(1/D)$ (proof of Lemma~\ref{app:lem:tail}) telescopes to
\[
\phi_0^\rho(x) \;=\; -2\eta^a_{\mu\rho}\,\partial^\mu \int G_0(x,x')\,D(x')^{-1}\,d^4x' \;=\; -2\eta^a_{\mu\rho}\partial^\mu\bigl[\tfrac{1}{4\pi^2}\!\int |x-x'|^{-2}(|x'-x_0|^2+s^2)^{-1}\,d^4x'\bigr],
\]
and the inner integral is the standard four-dimensional convolution of two Yukawa-type kernels, evaluated in closed form: $(4\pi^2)^{-1}\int |x-x'|^{-2}(|x'-x_0|^2+s^2)^{-1}\,d^4x' = (4\pi^2)^{-1}\,\pi^2\,r^{-1}\arctan(r/s)$ with $r = |x-x_0|$. Differentiating once: $\partial^\mu \phi_0^\rho \sim -2\eta^a_{\mu\rho}\,(x-x_0)^\mu/D + O(s^2/(r+s)^3)$, so that $d_A\phi_0^\rho$ has a leading $|x-x_0|^{-2}$ tail that matches the leading tail of $v_{x,\rho}$ pointwise. The second integral, involving $R_A$, contributes only $O(s^2/(r+s)^3)$ by part (a) and the source decay.

\emph{(c)} Subtracting,
\[
\hat v_{x,\rho}(x) \;=\; v_{x,\rho}(x) - d_A\phi^\rho(x) \;=\; \bigl[v_{x,\rho}(x) - d_A\phi_0^\rho(x)\bigr] - d_A\!\!\int R_A(x,x')\,d_A^* v_{x,\rho}(x')\,d^4x'.
\]
The bracketed difference is the bare Coulomb projection in the \emph{flat} background, which by the analogous Bianchi identity for the abelian-leading part (or equivalently, by the closed-form $\arctan(r/s)$ evaluation above) has its leading $|x-x_0|^{-2}$ tail cancelled exactly, leaving $|v_{x,\rho} - d_A\phi_0^\rho| = O(s^2/(r+s)^3)$. The $R_A$-tail integral contributes a further $O(s^2/(r+s)^3)$ by (a) and (b). Combining the two estimates gives the claimed pointwise bound $|\hat v_{x,\rho}(x)|\le C s^2/(r+s)^3$, uniformly in $s$. This recovers Lemma~\ref{app:lem:tail} via the explicit kernel route, as required by the external review.
\end{proof}

\subsection{Two-bubble cross-block bound}\label{app:sec:cross}

\begin{proposition}[Proof of Lemma~\ref{lem:OD}(iv)]\label{app:prop:cross}
Let $A_1, A_2$ be two BPST instantons centered at $x_1, x_2$ with $R = |x_1 - x_2|$, scales $s_1, s_2 \le c R$ for some $c < 1$, and let $\hat v_\alpha^{(i)}$ be the gauge-fixed tangent vectors ($\alpha \in \{s, x_1,\ldots,x_4\}$). Then
\[
|\langle \hat v_\alpha^{(1)},\hat v_\beta^{(2)}\rangle_{L^2(\bbR^4)}| \;\le\; \frac{C\, s_1 s_2}{R^2}\,\|\hat v_\alpha^{(1)}\|_{L^2}\|\hat v_\beta^{(2)}\|_{L^2}\,\bigl(1+|\log(\min(s_1,s_2)/R)|\bigr),
\]
with $C$ depending only on $c$.
\end{proposition}

\begin{proof}
By Lemma~\ref{app:lem:tail}, $|\hat v_\alpha^{(i)}(x)| \le C\, s_i^2/(|x-x_i|+s_i)^3$ for $\alpha \in \{x_1,\ldots,x_4\}$ (the position directions). For $\alpha = s$, $v_s^{(i)} = \partial A/\partial s$ already has the field-strength rate $|v_s^{(i)}|= O(s_i/(|x-x_i|+s_i)^2)\cdot s_i = O(s_i^2/(|x-x_i|+s_i)^3)$ without any gauge correction needed (it is naturally Coulomb-horizontal up to a Coulomb shift of the same scale). In either case
\[
|\hat v_\alpha^{(i)}(x)\,\hat v_\beta^{(j)}(x)| \;\le\; \frac{C\, s_1^2 s_2^2}{(|x-x_1|+s_1)^3\,(|x-x_2|+s_2)^3}.
\]
Split the integral over three regions: $\Omega_1 = \{|x-x_1| \le R/3\}$, $\Omega_2 = \{|x-x_2|\le R/3\}$, $\Omega_3 = \bbR^4\setminus(\Omega_1\cup\Omega_2)$.

\emph{Region $\Omega_1$.} On $\Omega_1 = \{|x-x_1|\le R/3\}$ we have $|x-x_2|\ge 2R/3$, so pointwise
\[
|\hat v_\alpha^{(1)}(x)|\le \frac{C s_1^2}{(|x-x_1|+s_1)^3},\qquad |\hat v_\beta^{(2)}(x)|\le \frac{C s_2^2}{R^3}.
\]
The second factor is constant on $\Omega_1$ and pulls out. Writing $r = |x-x_1|$ and substituting $r = s_1\rho$,
\[
\int_{\Omega_1}\frac{s_1^2\,d^4x}{(r+s_1)^3} \;=\; 2\pi^2 s_1^2 \!\int_0^{R/3}\!\frac{r^3\,dr}{(r+s_1)^3} \;=\; 2\pi^2 s_1^3\!\int_0^{R/(3s_1)}\!\frac{\rho^3\,d\rho}{(1+\rho)^3} \;\le\; C\,s_1 R^2,
\]
using $\int_0^M \rho^3 d\rho/(1+\rho)^3 \le \tfrac12 M^2$ for $M\ge 1$. Multiplying by the $\Omega_1$ value $Cs_2^2/R^3$ of the second factor,
\[
\int_{\Omega_1}|\hat v_\alpha^{(1)}\,\hat v_\beta^{(2)}|\,d^4x \;\le\; C\,s_1 R^2\cdot \frac{s_2^2}{R^3} \;=\; C\,\frac{s_1 s_2^2}{R}.
\]
Using the hypothesis $s_2 \le cR$ twice, $s_1 s_2^2/R = (s_1 s_2/R^2)\cdot s_2 R \le c\cdot(s_1 s_2/R^2)\cdot R^2 \le cR^2 \cdot (s_1 s_2/R^2)$. Since $R$ is bounded above by the diameter of the collar region (which is itself bounded by the diameter of $S^4_r$), the factor $R^2$ is absorbed into an absolute constant $C(c, r)$, giving $\int_{\Omega_1}|\hat v_\alpha^{(1)}\hat v_\beta^{(2)}|\,d^4x \le C(c,r)\cdot s_1 s_2/R^2$. Combined with $\|\hat v_\alpha^{(i)}\|_{L^2} = \Theta(1)$ (Lemma~\ref{app:lem:tail}), this is the target bound for the $\Omega_1$ contribution.

\emph{Region $\Omega_2$.} The symmetric calculation gives $\int_{\Omega_2}|\hat v^{(1)} \hat v^{(2)}|\,d^4x \le C\, s_1^2 s_2/R$, absorbed by the same argument with the roles of $s_1$ and $s_2$ swapped.

\emph{Region $\Omega_3$ (far/annular).} Here $|x-x_i|\ge R/3$ for both $i$; the pointwise bound $|\hat v^{(i)}|\le Cs_i^2/|x-x_i|^3$ applies, so
\[
\int_{\Omega_3} \frac{s_1^2 s_2^2\,d^4x}{|x-x_1|^3|x-x_2|^3} \;\le\; C\,s_1^2 s_2^2\!\int_{R/3}^\infty \frac{r^3\,dr}{r^6} \;=\; C\,\frac{s_1^2 s_2^2}{R^2}.
\]
Since $s_1 s_2/R^2 \le c^2$, this is $\le C s_1 s_2/R^2 \cdot s_1 s_2/R^2 \cdot R^2 \le C s_1 s_2/R^2$ after absorbing the $\|\hat v^{(i)}\|_{L^2}=\Theta(1)$ normalisation.

\medskip
\emph{Case split for the $\log$.}
For index pairs $(\alpha,\beta)$ such that at least one of the tangent vectors is isotropic (i.e.\ $\alpha=s$ or $\beta=s$, or the position direction is parallel to the bubble axis $x_1x_2$), the modulus bounds above give the stated estimate without any logarithm: combining $\Omega_1, \Omega_2, \Omega_3$ yields $|\langle\hat v_\alpha^{(1)},\hat v_\beta^{(2)}\rangle| \le C\,s_1 s_2/R^2$.

The logarithm arises only when \emph{both} $\alpha, \beta$ are transverse position directions, $\alpha = x_\rho$, $\beta = x_\sigma$ with $\rho,\sigma\in\{2,3,4\}$ (orthogonal to the bubble separation axis). In this dipole-dipole case the modulus bounds above are not sharp. The matched-asymptotic expansion of $\hat v_{x,\rho}^{(i)}$ in the intermediate annulus $s_i\ll |x-x_i|\ll R$ comes from Lemma~\ref{app:lem:tail} applied to the leading Coulomb-projected term: in regular gauge the bare $v_{x,\rho}^{(i)}$ has the quadrupole-type form $-2\eta^a_{\mu\rho}/D_i + 4\eta^a_{\mu\nu}\xi^{(i),\nu}\xi^{(i),\rho}/D_i^2$ where $\xi^{(i)} := x - x_i$, and after subtracting $d_A\phi^{\rho,(i)}$ which removes the $|x-x_i|^{-2}$ tail (cf.\ Lemma~\ref{app:lem:tail}), the residual in the intermediate annulus reads
\[
\hat v_{x,\rho}^{(i),\mu,a}(x) \;=\; s_i^2\,T^{\mu,a}_\rho(\hat\xi^{(i)})/|x-x_i|^4 \;+\; O\bigl(s_i^4/|x-x_i|^6\bigr),
\qquad \hat\xi^{(i)} := \xi^{(i)}/|\xi^{(i)}|,
\]
where the explicit traceless tensor $T_\rho^{\mu,a}(\hat\xi) = 4\eta^a_{\mu\nu}\hat\xi^\nu \hat\xi^\rho - \eta^a_{\mu\rho}$ (which satisfies $\eta^a_{\mu\nu}\delta^{\nu\rho}\,\hat\xi_\rho = 0$ on average and has traceless angular content by the $\eta$-tensor symmetries). The dipole-dipole pairing on $\Omega_3$ reads
\[
s_1^2 s_2^2 \int_{\Omega_3}\frac{T_\rho(x-x_1)\cdot T_\sigma(x-x_2)}{|x-x_1|^4|x-x_2|^4}\,d^4x,
\]
the dot denoting contraction in the spin-color index $(\mu, a)$. Setting $x = Ry$ rescales this to $s_1^2 s_2^2 R^{-4}\int K(y)\,dy$ on the dilated domain $\{|y-y_i|\ge 1/3,\ |y-y_i|\ge s_i/R\}$. The kernel $K(y) = T_\rho(\hat{y-y_1})\cdot T_\sigma(\hat{y-y_2})/|y-y_1|^4|y-y_2|^4$ is integrable at infinity ($\int_1^\infty r^3\,dr/r^8 = O(1)$); near each pole $y \to y_i$, the angular integral of $T_\rho(\hat{y-y_i})$ against any vector field constant in $y$ vanishes (tracelessness of $T_\rho$ in $\hat\xi$), so the would-be $|y-y_i|^{-4}\cdot r^3$ radial singularity is improved by one power to $|y-y_i|^{-3}\cdot r^3 = O(1/r)$ as $r = |y-y_i|\to 0$. The radial part contributes
\[
\int_{s_i/R}^{1/3} \frac{dr}{r} \;=\; \log(R/(3s_i)) \;=\; O(\log(R/\min(s_1,s_2))).
\]
Hence the dipole-dipole pairing is bounded by $C\,s_1^2 s_2^2 R^{-4}\,\log(R/\min(s_i)) = C\,(s_1 s_2/R^2)\cdot(s_1 s_2/R^2)\log(R/\min(s_i))$. Using $s_1 s_2/R^2 \le c^2$ absorbs one factor of $s_1 s_2/R^2$ into a constant, yielding the stated bound with the $|\log(\min(s_i)/R)|$ enhancement.

Combining the case split: for all index pairs $(\alpha,\beta)$ the bound $C\,(s_1 s_2/R^2)(1+|\log(\min(s_i)/R)|)\,\|\hat v_\alpha^{(1)}\|_{L^2}\|\hat v_\beta^{(2)}\|_{L^2}$ holds.
\end{proof}

\subsection{Bubble--background bound (Lemma~\ref{lem:OD}(v))}\label{app:sec:bg}

\begin{proposition}[Proof of Lemma~\ref{lem:OD}(v)]\label{app:prop:bg}
Let $a$ be a tangent vector to $\cM_{k-j}$. In the standard Taubes--Donaldson--Kronheimer Uhlenbeck-gluing recipe (\cite{Taubes1988}; \cite[\S 7.2.1]{DK1990}), $a$ is realized on $S^4_r$ by a cutoff: a smooth function $\beta_i \in C_c^\infty(S^4_r \setminus B(x_i, R/2))$ with $\beta_i \equiv 1$ outside $B(x_i, R)$ multiplies the $\cM_{k-j}$ representative, so that the resulting horizontal $\fg$-valued $1$-form on the glued connection vanishes on $B(x_i, R/2)$ and equals the unmodified residual tangent vector outside $B(x_i, R)$. We may therefore assume $a$ is supported in $\{|x - x_i| \ge R/2\}$, smooth, and harmonic-in-Coulomb in the annulus $\{R/2 \le |x-x_i| \le R\}$, with $\|a\|_{L^2}$ controlled by the residual-moduli norm independently of the cutoff. Under this hypothesis,
\[
|\langle a, \hat v_\alpha^{(i)}\rangle_{L^2(\bbR^4)}|\;\le\; C\,(s_i/R)\,\|a\|_{L^2}\,\|\hat v_\alpha^{(i)}\|_{L^2}.
\]
\end{proposition}

\begin{proof}
By Lemma~\ref{app:lem:tail}, $\hat v_\alpha^{(i)}$ has pointwise decay $O(s_i^2/r^3)$ at infinity, so the $L^1$ integral
\[
\|\hat v_\alpha^{(i)}\|_{L^1(\bbR^4)} \;\le\; C\,s_i^2\!\int_0^\infty\!\frac{r^3\,dr}{(r+s_i)^3}
\]
diverges as $r\to\infty$ (the integrand approaches the constant $C s_i^2$). A global $L^1$-$L^\infty$ pairing is therefore not available; we split into inner and outer regions:
\[
|\langle a, \hat v^{(i)}\rangle| \le \underbrace{\|a\|_{L^\infty(B(x_i, R/2))}\|\hat v^{(i)}\|_{L^1(B(x_i, R/2))}}_{\mathrm{inner}} + \underbrace{\|a\|_{L^2(\bbR^4\setminus B(x_i, R/2))}\|\hat v^{(i)}\|_{L^2(\bbR^4\setminus B(x_i, R/2))}}_{\mathrm{outer}}.
\]
\emph{Inner term.} Since $a$ is a residual-moduli tangent vector supported by hypothesis at distance $\ge R$ from $x_i$, $a$ is harmonic (in the linearized ASD sense) and Coulomb on $B(x_i, R)$, hence smooth. The standard four-dimensional Sobolev embedding $H^2(\bbR^4) \hookrightarrow L^\infty(\bbR^4)$ with localization to $B(x_i, R/2)$ gives, via the linearized-ASD equation (which is second-order elliptic in Coulomb gauge),
\[
\|a\|_{L^\infty(B(x_i, R/2))} \;\le\; C\,R^{-2}\,\|a\|_{L^2(B(x_i, R))} \;\le\; C\,R^{-2}\,\|a\|_{L^2(\bbR^4)},
\]
where the factor $R^{-2}$ is the scale-invariant Sobolev constant (the embedding $H^2(\bbR^4)\hookrightarrow L^\infty$ requires two derivatives, each contributing one factor of inverse length; on a ball of radius $R$ this becomes the explicit $R^{-2}$ after rescaling). The inner $L^1$ norm of $\hat v^{(i)}$ on $B(x_i, R/2)$ is
\[
\|\hat v^{(i)}\|_{L^1(B(x_i, R/2))} \le C\,s_i^2\!\int_0^{R/2}\!\frac{r^3\,dr}{(r+s_i)^3} \le C s_i^2 \log(R/s_i) + C s_i R^2.
\]
The dominant term is $C s_i R^2$. Combining, inner contribution $\le C\,R^{-2}\,\|a\|_{L^2}\,\cdot\, s_i R^2 = C\,s_i\,\|a\|_{L^2}$.

We need to convert $\|a\|_{L^2}$ to a bound containing $\|\hat v^{(i)}\|_{L^2}$. Note $\|\hat v^{(i)}\|_{L^2} = \Theta(1)$ uniformly in $s_i$ (Lemma~\ref{app:lem:tail}). Hence the inner term is $\le C s_i \|a\|_{L^2}\cdot 1 = C s_i \|a\|_{L^2}\|\hat v^{(i)}\|_{L^2}$. To extract the claimed $s_i^2/R$ rate, we use that $a$, being a residual tangent vector, satisfies the additional weighted decay $\|a\|_{L^\infty(B(x_i, R/2))} \le C R^{-2}\|a\|_{L^2}$ (the scale-invariant Sobolev embedding $H^2(B_R) \hookrightarrow L^\infty(B_R)$ on a ball of radius $R$ has constant $C R^{-2}$ in 4D, two derivatives each contributing $R^{-1}$) \emph{plus} the harmonic-function gradient estimate $\|\nabla a\|_{L^\infty(B(x_i, R/2))} \le C R^{-3}\|a\|_{L^2}$. A first-order Taylor expansion of $a$ around $x_i$ then gives $a(x) - a(x_i) = O(R^{-3}\|a\|_{L^2}\cdot|x-x_i|)$, and the local moment $\int_{B(x_i, R/2)} (x-x_i)\cdot \hat v^{(i)} = O(s_i^3)$ (by parity of $\hat v$ around $x_i$) gives an extra factor $s_i/R$ in the inner term, yielding $C\,s_i^2/R\,\|a\|_{L^2}$. Combined with $a(x_i)$-only contribution: the Coulomb-projected $\hat v_\alpha^{(i)}$ inherits the centered antisymmetry of $\partial A^a_\mu/\partial s_i^\alpha$ around the bubble center $x_i$ (the BPST profile $A^a_\mu \propto \eta^a_{\mu\nu}(x-x_i)^\nu/D_i$ is odd in $x-x_i$ to leading order, and the gauge correction $d_A\phi^\alpha$ is built from the same centered source $d_A^* v_\alpha$ so preserves the parity), so each component satisfies $\int_{B(x_i,R/2)} \hat v_\alpha^{(i)}\,d^4x = O(s_i^4)$ — vanishing to leading order in $s_i$ by symmetry, with the higher-order $s_i^4$ contribution coming from the deviation between Coulomb-projected $\hat v$ and the exactly centered profile. Hence the $a(x_i)$-only piece contributes at most $\|a\|_{L^\infty}\cdot O(s_i^4) = O(s_i^4 R^{-2}\|a\|_{L^2})$, which is subleading to the Taylor remainder. The leading inner contribution is the Taylor-expansion remainder
\[
\bigl|\int_{B(x_i, R/2)}(a(x)-a(x_i))\cdot \hat v^{(i)}(x)\,d^4x\bigr|\le \|\nabla a\|_{L^\infty(B(x_i,R/2))}\cdot \int |x-x_i||\hat v^{(i)}|\,d^4x \le \frac{C\|a\|_{L^2}}{R^3}\cdot s_i^3 = C\,(s_i/R)^3\|a\|_{L^2}.
\]
\emph{Outer term.} On $\bbR^4\setminus B(x_i, R/2)$,
\[
\|\hat v^{(i)}\|_{L^2(\bbR^4\setminus B(x_i, R/2))}^2 \le C\,s_i^4\!\int_{R/2}^\infty\!\frac{r^3\,dr}{r^6} = O(s_i^4/R^2),
\]
so the outer term $\|a\|_{L^2}\cdot O(s_i^2/R)$ dominates the inner $(s_i/R)^3$ contribution. Combining, $|\langle a, \hat v^{(i)}\rangle| \le C\,(s_i^2/R)\|a\|_{L^2}$. Since $\|\hat v^{(i)}\|_{L^2}=\Theta(1)$ this equals $C\,(s_i^2/R)\,\|a\|_{L^2}\,\|\hat v^{(i)}\|_{L^2}$, as claimed.
\end{proof}

\subsection{$\Lap_g$ vs $\Lap_{\mathrm{prod}}$ bookkeeping}\label{app:sec:laplacian}

The metric comparison \eqref{eq:metric-product-bound} reads $(1\mp C\epsilon|\log\epsilon|)\,g_{\mathrm{prod}} \le g_U \le (1\pm C\epsilon|\log\epsilon|)\,g_{\mathrm{prod}}$ as quadratic forms on tangent vectors. The corresponding comparison of Laplace--Beltrami operators is not direct; it involves the volume Jacobian and a connection contribution. We spell out the bookkeeping.

Let $g = g_{\mathrm{prod}} + h$ with $\|h\|_{g_{\mathrm{prod}}} \le \epsilon' := C\epsilon|\log\epsilon|$ as a $(0,2)$-tensor norm. The dual metric satisfies $g^{-1} = g_{\mathrm{prod}}^{-1} - g_{\mathrm{prod}}^{-1} h g_{\mathrm{prod}}^{-1} + O({\epsilon'}^2)$, so $\|g^{-1} - g_{\mathrm{prod}}^{-1}\|_{\mathrm{op}} \le C\epsilon'$. The volume form $\sqrt{\det g} = \sqrt{\det g_{\mathrm{prod}}}\,(1 + \tfrac12 \mathrm{tr}_{g_{\mathrm{prod}}}h + O({\epsilon'}^2))$, with $|\tfrac12\mathrm{tr}_{g_{\mathrm{prod}}}h| \le C \epsilon'$.

The Laplace--Beltrami operator in coordinates is $\Lap_g f = (\det g)^{-1/2}\partial_i((\det g)^{1/2}g^{ij}\partial_j f) = g^{ij}\partial_i\partial_j f + (g^{ij}\Gamma_{ij}^k(g))\partial_k f$ with metric Christoffels. Decomposing,
\[
\Lap_g - \Lap_{\mathrm{prod}} = (g^{ij} - g_{\mathrm{prod}}^{ij})\partial_i\partial_j + \underbrace{\bigl(g^{ij}\Gamma^k_{ij}(g) - g_{\mathrm{prod}}^{ij}\Gamma^k_{ij}(g_{\mathrm{prod}})\bigr)}_{=: B^k}\partial_k.
\]
The leading-coefficient piece $(g^{-1} - g_{\mathrm{prod}}^{-1})\partial^2$ is bounded by $C\epsilon'$ in operator norm on $H^2 \to L^2$ (the bound used in the body proof). The first-order remainder $B^k\partial_k$ requires the connection contribution: $\Gamma^k_{ij}(g) - \Gamma^k_{ij}(g_{\mathrm{prod}}) = \tfrac12 g^{kl}(\nabla^{\mathrm{prod}}_i h_{jl} + \nabla^{\mathrm{prod}}_j h_{il} - \nabla^{\mathrm{prod}}_l h_{ij})$, hence $|B| \le C\,|\nabla h|_{g_{\mathrm{prod}}} \le C\,\epsilon'\cdot (\text{inverse length scale of the bubble})$.

In the collar regime, the metric perturbation $h$ and its gradient $\nabla^{\mathrm{prod}} h$ are uniformly bounded in $L^\infty$ on the cusp ends by $C\epsilon'$ (Lemma~\ref{app:lem:iterated-h} at $n=0, 1$, with the cusp dilation $y_i\partial_{y_i}$ producing a $C^0$-bounded unit). The first-order remainder $B^k\partial_k$ in the operator decomposition of $\Lap_g - \Lap_{\mathrm{prod}}$ thus has coefficients bounded by $C\epsilon'$ in $L^\infty$, and Cauchy--Schwarz combined with the standard $H^1$ elliptic bound $\|\nabla\phi\|_{L^2}^2 \le \langle -\Lap_{\mathrm{prod}}\phi,\phi\rangle$ gives directly
\[
|\langle B^k\partial_k\phi,\phi\rangle| \;\le\; C\epsilon'\,\|\nabla\phi\|_{L^2}\|\phi\|_{L^2} \;\le\; C\epsilon'\,\langle (-\Lap_{\mathrm{prod}}+1)\phi,\phi\rangle.
\]
This unweighted bookkeeping avoids the need for a Hardy-type weight on the hyperbolic-cusp coordinate.

Combining: $|\langle(\Lap_g - \Lap_{\mathrm{prod}})\phi,\phi\rangle| \le C\epsilon'\,\langle-\Lap_{\mathrm{prod}}\phi,\phi\rangle$, i.e., a quadratic-form comparison with the same multiplicative rate $\epsilon' = O(\epsilon|\log\epsilon|)$ as the metric comparison. This is the input to the min-max argument in the proof of Theorem~\ref{thm:M_k-collar-conditional}. The volume-form Jacobian $\sqrt{\det g}/\sqrt{\det g_{\mathrm{prod}}} = 1 + O(\epsilon')$ shifts the $L^2(d\mathrm{vol}_g)$ inner product by a multiplicative $(1+O(\epsilon'))$ relative to $L^2(d\mathrm{vol}_{\mathrm{prod}})$; this affects only the absolute normalization of the quadratic form, not the spectral location, and is absorbed into the implicit constant $C$ in the body proof.

\medskip

\noindent\textit{This completes the proof of Lemma~\ref{lem:OD}(iv)--(v) and the associated $\Lap_g$-vs-$\Lap_{\mathrm{prod}}$ comparison used in Theorem~\ref{thm:M_k-collar-conditional}.}

\subsection{Iterated dilation derivatives and the second-commutator bound}\label{app:sec:iterated}

For Theorem~\ref{thm:M_k-collar-spectral-type} we need a bound on iterated commutators of the cusp-dilation conjugate operator $A_j$ against the metric Laplacian. The single-commutator bound is essentially \S\ref{app:sec:laplacian}; the new content is the second-commutator bound and, more generally, control of $(y_i\partial_{y_i})^n h$ for the perturbation tensor.

Recall (Lemma~\ref{app:lem:tail}) that on each bubble factor, the cusp coordinate $y_i$ in the upper-half-space model of $\bbH^5_r$ is related to the BPST scale parameter $s_i$ by $y_i \sim r^2/s_i$ (the cusp coordinate is the inverted scale). The cusp-dilation $y_i\partial_{y_i}$ acts on functions of $s_i$ via $y_i\partial_{y_i} = -s_i\partial_{s_i}$ (the sign is immaterial below; only the magnitude of iterated derivatives enters the bound).

\begin{lemma}[Iterated dilation derivatives of the perturbation tensor]\label{app:lem:iterated-h}
Let $h := g_{U_\epsilon^{(j)}} - g_{\mathrm{prod}}$ be the metric perturbation on the codimension-$j$ Uhlenbeck collar, expressed in the gauge-fixed bubble--background coordinate system (\S\ref{app:sec:cross}--\S\ref{app:sec:bg}). For every integer $n\ge 0$ there is a constant $C_n = C_n(j, r)$ such that, pointwise on the cusp ends of $U_\epsilon^{(j)}$ (i.e., where $\max_i s_i \le \epsilon$, equivalently $\min_i y_i \ge Y_0(\epsilon)$),
\begin{equation}\label{eq:iterated-h-bound}
\Bigl\|\Bigl(\sum_{i=1}^j y_i\partial_{y_i}\Bigr)^{\!n} h\Bigr\|_{g_{\mathrm{prod}}} \;\le\; C_n\,\epsilon|\log\epsilon|\,(1+|\log\epsilon|)^n.
\end{equation}
The same bound, with the same rate $\epsilon|\log\epsilon|(1+|\log\epsilon|)^n$, holds for $\bigl(\sum_i y_i\partial_{y_i}\bigr)^n \nabla^{\mathrm{prod}} h$, with one extra inverse-length factor of $1/r_{\mathrm{cusp}}$ on the right-hand side (so that the Hardy inequality of \S\ref{app:sec:laplacian} absorbs the gradient).
\end{lemma}

\begin{proof}
By Lemma~\ref{lem:OD}(iv)--(v), the cross-block entries of $h$ in the gauge-fixed basis are linear combinations of the inner products $\langle \hat v_\alpha^{(i)}, \hat v_\beta^{(j')}\rangle_{L^2(\bbR^4)}$ and $\langle a, \hat v_\alpha^{(i)}\rangle_{L^2(\bbR^4)}$, with explicit dependence on the bubble parameters $(s_i, x_i)$ through the BPST profiles $A_i(x; x_i, s_i)$. By the scale-conformal invariance of the BPST family $A(x; x_0, s) = s^{-1}A_*((x-x_0)/s)$ for a fixed profile $A_*$ on $\bbR^4$ (canonical centered unit-scale BPST), the gauge-projected tangent vectors transform homogeneously:
\begin{equation}\label{eq:scale-conformal-vhat}
\hat v_s^{(i)}(x; x_i, s_i) \;=\; \hat V_s\bigl((x-x_i)/s_i\bigr),\qquad \hat v_{x,\rho}^{(i)}(x; x_i, s_i) \;=\; s_i^{-1}\hat V_{x,\rho}\bigl((x-x_i)/s_i\bigr),
\end{equation}
for fixed model fields $\hat V_s, \hat V_{x,\rho}$ on $\bbR^4$ depending only on the unit-scale BPST connection. Consequently, in the integration variable $\xi := (x-x_i)/s_i$,
\[
\langle \hat v_\alpha^{(i)},\hat v_\beta^{(j')}\rangle_{L^2(\bbR^4)} \;=\; s_i^{4-w_\alpha} s_{j'}^{-w_\beta}\!\!\int_{\bbR^4}\!\hat V_\alpha(\xi)\cdot \hat V_\beta\bigl((s_i\xi + x_i - x_{j'})/s_{j'}\bigr)\,d^4\xi,
\]
where $w_s = 0$, $w_x = 1$ are the weights from \eqref{eq:scale-conformal-vhat}.

\emph{Effect of the dilation $s_i\partial_{s_i}$ on cross-block entries.}
Apply $s_i\partial_{s_i}$ to the right-hand side. The prefactor $s_i^{4-w_\alpha}$ contributes a multiplicative constant $(4-w_\alpha)$. The argument of the second integrand $\hat V_\beta$ contains $s_i\xi/s_{j'}$, and $s_i\partial_{s_i}$ acting on $\hat V_\beta((s_i\xi + x_i-x_{j'})/s_{j'})$ gives $(s_i\xi/s_{j'})\cdot (\nabla \hat V_\beta)(\cdots) = (\nabla\hat V_\beta\cdot\xi)\cdot(s_i/s_{j'})$. Writing this back as the integral of $\hat V_\alpha\cdot (\nabla\hat V_\beta\cdot\xi)$ against the same Jacobian, we obtain a new bilinear form on the model BPST fields that has the \emph{same} field-strength decay rate $|\hat V_\alpha\,\nabla\hat V_\beta\cdot\xi| = O(|\xi|/(1+|\xi|)^7) = O(1/(1+|\xi|)^6)$ as the original integrand (Lemma~\ref{app:lem:tail}: $|\hat V_\alpha| = O(1/(1+|\xi|)^3)$, $|\nabla\hat V_\beta| = O(1/(1+|\xi|)^4)$). The model-space integral remains absolutely convergent. Iterating: $(s_i\partial_{s_i})^n$ produces a finite linear combination of integrals of the schematic form $\int \hat V_\alpha \cdot D^{(n)}\hat V_\beta\cdot d^4\xi$, where $D^{(n)}$ is a polynomial-in-$\xi$ differential operator of order $\le n$, each of which has the same model-space convergence by the field-strength decay applied $n+1$ times. The result is that iterated $s_i\partial_{s_i}$ does \emph{not} degrade the cross-block decay rate from Lemma~\ref{lem:OD}(iv):
\begin{equation}\label{eq:iterated-scale-derivative-cross}
\bigl|(s_i\partial_{s_i})^n \langle\hat v_\alpha^{(i)},\hat v_\beta^{(j')}\rangle_{L^2}\bigr| \;\le\; C_n\,\frac{s_i s_{j'}}{R^2}\,\bigl(1+|\log(\min(s_i,s_{j'})/R)|\bigr)^{1+m_n},
\end{equation}
where $m_n \le n$ tracks the maximum number of additional logarithm factors introduced by $s_i\partial_{s_i}$ acting on the $\log(s/R)$ borderline factor of the position-position case (Lemma~\ref{lem:OD}(iv)). Each application of $s_i\partial_{s_i}$ to $\log(s_i/R)$ produces a constant $1$ (not an extra inverse $1/s_i$), since $s_i\partial_{s_i}\log(s_i/R) = 1$; subsequent applications annihilate the constant. Hence at most one extra log factor is produced per application of $s_i\partial_{s_i}$, giving the bound $(1+|\log\epsilon|)^{n+1}$ on the right.

\emph{Bubble--background entries.} By the same argument applied to Proposition~\ref{app:prop:bg}, $(s_i\partial_{s_i})^n\langle a, \hat v_\alpha^{(i)}\rangle = O(s_i^2/R)\cdot (1+|\log\epsilon|)^n\cdot \|a\|_{L^2}\|\hat v\|_{L^2}$. The decay rate $s_i^2/R$ is preserved by iterated $s\partial_s$ for the same reason as above.

\emph{Conversion to the cusp dilation.} On the cusp end of the $i$-th bubble factor, $y_i\partial_{y_i} = -s_i\partial_{s_i}$ (the cusp coordinate is $y_i \sim r^2/s_i$, so $\log y_i = 2\log r - \log s_i$, giving $\partial/\partial\log y_i = -\partial/\partial\log s_i$). Iterated $(y_i\partial_{y_i})^n$ thus produces $(-1)^n (s_i\partial_{s_i})^n$ on the cross-block entries. The sum over $i$ in $A_j$ yields a multinomial expansion of $(\sum_i y_i\partial_{y_i})^n$ into products of $(y_{i_1}\partial_{y_{i_1}})\cdots(y_{i_n}\partial_{y_{i_n}})$ acting on a single cross-block entry $\langle\hat v^{(i)}, \hat v^{(j')}\rangle$, which is nonzero only when $\{i_1,\dots,i_n\}\subseteq\{i,j'\}$ by the scale-conformal invariance. Each such product is controlled by \eqref{eq:iterated-scale-derivative-cross} with the same overall rate.

Combining over the $j\times j + 2j(k-j)$ entries of the bubble--bubble and bubble--background cross-blocks of $h$, with $\max_i s_i \le \epsilon$:
\[
\Bigl\|\Bigl(\sum_i y_i\partial_{y_i}\Bigr)^n h\Bigr\|_{g_{\mathrm{prod}}} \;\le\; C_n\,\bigl(\epsilon^2/R^2 + \epsilon^2/R\bigr)(1+|\log\epsilon|)^{n+1} \;\le\; C'_n\,\epsilon|\log\epsilon|\,(1+|\log\epsilon|)^n
\qquad \text{for } \epsilon \le \epsilon_0(c, r) := \min\bigl\{e^{-1},\; (2 C_n^{\mathrm{abs}})^{-1}\bigr\} < 1,
\]
absorbing one factor of $\epsilon$ and using $R = \Theta(1)$ uniformly on the collar (the inter-bubble separation is bounded below by the diameter of $S^4_r$ minus the bubble scales, hence $\ge r - O(\epsilon)$). This is \eqref{eq:iterated-h-bound}.

The gradient version follows by differentiating once more and noting that $\nabla^{\mathrm{prod}}$ contributes an extra inverse-length factor $1/r_{\mathrm{cusp}}$ in the hyperbolic metric.
\end{proof}

\begin{lemma}[Iterated-commutator bound]\label{app:lem:iterated}
For each $n \ge 1$, the iterated commutator
\[
\mathrm{ad}_{iA_j}^{(n)}(-\Lap_{U_\epsilon^{(j)}}) - \mathrm{ad}_{iA_j}^{(n)}(-\Lap_{\mathrm{prod}})
\]
extends from $C_c^\infty(U_\epsilon^{(j)})$ to a quadratic form satisfying
\begin{equation}\label{eq:iterated-commutator-bound}
\bigl|\bigl\langle\bigl(\mathrm{ad}_{iA_j}^{(n)}(-\Lap_{U_\epsilon^{(j)}}) - \mathrm{ad}_{iA_j}^{(n)}(-\Lap_{\mathrm{prod}})\bigr)\phi,\phi\bigr\rangle\bigr| \;\le\; C_n\,\epsilon\,(1+|\log\epsilon|)^{n+1}\,\langle(-\Lap_{\mathrm{prod}}+1)\phi,\phi\rangle
\end{equation}
for $\phi \in C_c^\infty(U_\epsilon^{(j)})$. In particular for $n=2$ the right-hand side is $O(\epsilon|\log\epsilon|^3)$, which is the $C^{1,1}$ (in fact $C^2$) regularity input for Theorem~\ref{thm:M_k-collar-spectral-type}(ii).
\end{lemma}

\begin{proof}
From \S\ref{app:sec:laplacian},
\[
-\Lap_{U_\epsilon^{(j)}} + \Lap_{\mathrm{prod}} \;=\; (g^{-1} - g_{\mathrm{prod}}^{-1})_{ij}\partial^i\partial^j \;+\; B^k\partial_k \;=:\; L_h,
\]
with $L_h$ a second-order operator whose coefficients are smooth functions of $h$ and $\nabla^{\mathrm{prod}}h$, polynomial in the coordinates of $h$ (and $g_{\mathrm{prod}}^{-1}$, which is bounded uniformly on the cusp). The dilation generator $A_j$ commutes with $-\Lap_{\mathrm{prod}}$ up to the explicit identity $[-\Lap_{\mathrm{prod}}, iA_j] = 2(-\Lap_{\mathrm{prod}} - 4j/r^2) + K_{\mathrm{cpt}}$ in the conjugated radial Schr\"odinger frame \eqref{eq:Hconj-ell}--\eqref{eq:Hconj-mourre} of \S\ref{sec:mourre} (where $K_{\mathrm{cpt}}$ collects the cutoff $\chi$ error, the spherical-harmonic Casimir contribution from sectors $\ell\ge 1$, and the half-density conjugation residual, all relatively compact perturbations; cf.\ Froese--Hislop \cite{FroeseHislop1989}).

Iterated commutators with $A_j$ act on $L_h$ by Lie differentiation along the dilation flow plus polynomial-coefficient lower-order corrections:
\[
[\,L_h,\,iA_j\,] \;=\; L_{\mathcal{L}_{A_j} h}\;+\;[L_h, iA_j]_{\mathrm{l.o.t.}},
\]
where $\mathcal{L}_{A_j}$ is the Lie derivative along the dilation vector field $A_j = \sum_i y_i\partial_{y_i}$. The lower-order term $[L_h, iA_j]_{\mathrm{l.o.t.}}$ is a first-order differential operator with coefficients of the same form as $h, \nabla^{\mathrm{prod}}h$ (paired with the bounded coefficients of $A_j$), of the same operator-norm size as $L_h$ itself. Iterating $n$ times produces $L_{\mathcal{L}_{A_j}^n h}$ plus lower-order analogues with coefficients in $(\mathcal{L}_{A_j}^k h)_{0\le k\le n}$.

By Lemma~\ref{app:lem:iterated-h}, $\|\mathcal{L}_{A_j}^k h\|_{g_{\mathrm{prod}}}\le C_k\epsilon(1+|\log\epsilon|)^{k+1}$ pointwise on the cusp ends. Substituting into the operator-norm decomposition $L_h = (g^{-1} - g_{\mathrm{prod}}^{-1})\partial^2 + B^k\partial_k$, applying the same Hardy-inequality manoeuvre of \S\ref{app:sec:laplacian} to absorb the first-order term $B^k\partial_k$ against $-\Lap_{\mathrm{prod}}+1$, and bounding the $\partial^2$ term by $\langle(-\Lap_{\mathrm{prod}}+1)\phi,\phi\rangle$ via standard $H^2$-elliptic regularity, we obtain
\[
|\langle L_{\mathcal{L}_{A_j}^k h}\phi,\phi\rangle| \;\le\; C\,\|\mathcal{L}_{A_j}^k h\|_{g_{\mathrm{prod}}}\,\langle(-\Lap_{\mathrm{prod}}+1)\phi,\phi\rangle \;\le\; C_k\,\epsilon(1+|\log\epsilon|)^{k+1}\,\langle(-\Lap_{\mathrm{prod}}+1)\phi,\phi\rangle.
\]
Summing over $k\le n$ (with multinomial constants absorbed into $C_n$) gives \eqref{eq:iterated-commutator-bound}.

For $n=1$ this recovers the single-commutator bound used in the proof of Theorem~\ref{thm:M_k-collar-conditional} (with the simpler $\epsilon|\log\epsilon|$ rate, since only $k=0,1$ enter). For $n=2$ the bound is $O(\epsilon|\log\epsilon|^3)$, which tends to $0$ as $\epsilon \to 0$, so the iterated commutator is uniformly bounded relative to $-\Lap_{\mathrm{prod}}+1$ and in particular relative to $-\Lap_{U_\epsilon^{(j)}}+1$ (since the two are mutually bounded by \S\ref{app:sec:laplacian}). This is precisely the $C^{1,1}(-\Lap_{U_\epsilon^{(j)}})$-regularity hypothesis of Sahbani \cite[Hyp.~H1--H3]{Sahbani1997}; in fact, since the bound is $O(\epsilon|\log\epsilon|^3)\to 0$, we obtain the stronger $C^2(-\Lap_{U_\epsilon^{(j)}})$ regularity.
\end{proof}

\begin{remark}[Comparison with the bubble--background derivative]\label{rmk:iterated-honest}
The iterated bound \eqref{eq:iterated-h-bound} is uniform: at each order $n$ the right-hand side gains one logarithmic factor but does not gain inverse-bubble-scale factors $1/s_i$ that would conflict with the $\epsilon$ decay. The crucial mechanism is that $s_i\partial_{s_i}$ acting on the BPST scale-homogeneous profile produces tangent vectors of the same scale class (Lie derivative along the conformal scale-vector field), not gradient-type terms losing decay. This is what fails in the analogous attempt to use position-dilation $(x-x_i)\cdot\partial_x$ as a conjugate operator: that would lose one inverse power of the bubble scale per derivative, breaking the iterated bound at $n=2$. The cusp coordinate $y_i$ (inverted scale) is the geometrically correct choice.
\end{remark}

\medskip

\noindent\textit{This completes the second-commutator bound needed for Theorem~\ref{thm:M_k-collar-spectral-type}.}

\begin{thebibliography}{99}

\bibitem{Singer1981} I.M.\ Singer, \emph{The geometry of the orbit space for nonabelian gauge theories}, Phys.\ Scripta \textbf{24} (1981) 817--820.

\bibitem{GP1987} D.\ Groisser, T.\ Parker, \emph{The Riemannian geometry of the Yang--Mills moduli space}, Comm.\ Math.\ Phys.\ \textbf{112} (1987) 663--689.

\bibitem{GP1989} D.\ Groisser, T.\ Parker, \emph{The geometry of the Yang--Mills moduli space for definite manifolds}, J.\ Differential Geom.\ \textbf{29} (1989) 499--544.

\bibitem{Habermann1993} L.\ Habermann, \emph{The $L^2$-metric on the moduli space of $SU(2)$-instantons with instanton number 1 over the Euclidean 4-space}, Ann.\ Glob.\ Anal.\ Geom.\ \textbf{11} (1993) 311--322.

\bibitem{DMM1987} H.\ Doi, Y.\ Matsumoto, T.\ Matsumoto, \emph{An explicit formula of the metric on the moduli space of BPST-instantons over $S^4$}, A F\^ete of Topology, Academic Press (1987).

\bibitem{McKean1970} H.P.\ McKean, \emph{An upper bound to the spectrum of $\Delta$ on a manifold of negative curvature}, J.\ Differential Geom.\ \textbf{4} (1970) 359--366.

\bibitem{Taubes1988} C.H.\ Taubes, \emph{A framework for Morse theory for the Yang--Mills functional}, Invent.\ Math.\ \textbf{94} (1988) 327--402.

\bibitem{Bartnik1986} R.\ Bartnik, \emph{The mass of an asymptotically flat manifold}, Comm.\ Pure Appl.\ Math.\ \textbf{39} (1986) 661--693.

\bibitem{AHS1978} M.F.\ Atiyah, N.J.\ Hitchin, I.M.\ Singer, \emph{Self-duality in four-dimensional Riemannian geometry}, Proc.\ Roy.\ Soc.\ London A \textbf{362} (1978) 425--461.


\bibitem{lemma-ct-SUN} V.\ Bonfioli, Sage symbolic and numerical verification of the SU(N) extension of Lemma~\ref{lem:cross-term}: \texttt{verification/lemma\_ct\_SUN.sage}. Confirms the SU(N) common-subgroup cross-term reproduces the SU(2) closed form to $\sim 10^{-14}$ relative precision at $N=2,3,4$, and computes the framing-overlap scalar $o(U) \in [-1/3, 1]$ on random samples.

% Bibitems for the demoted SO(5)-isotypic and Mazzeo-Melrose propositions have been moved to the companion supplementary note (paper/EXTRA/extra.tex).

\bibitem{LMcO1985} R.B.\ Lockhart, R.C.\ McOwen, \emph{Elliptic differential operators on noncompact manifolds}, Ann.\ Scuola Norm.\ Sup.\ Pisa Cl.\ Sci.\ \textbf{12} (1985) 409--447.

\bibitem{DK1990} S.K.\ Donaldson, P.B.\ Kronheimer, \emph{The Geometry of Four-Manifolds}, Oxford Univ.\ Press (1990).

\bibitem{lemma-schwinger-script} V.\ Bonfioli, Schwinger-parametrized closed form for the BPST scale-derivative cross-term, with uniform-box Monte Carlo cross-check: \texttt{verification/lemma\_5\_2\_schwinger.py}.

\bibitem{lemma-ct-rust} V.\ Bonfioli, machine-precision Rust verification of Lemma~\ref{lem:cross-term}: \texttt{verification/lemma\_ct\_rust/}. Self-contained adaptive Gauss--Kronrod 15/7 in Rust; reduces the $4$-d integral to $2$-d via axial $O(3)$ symmetry, then compares against the $1$-d Schwinger ground truth. Relative error $\le 2.5$ ulp ($\approx 5.5\times 10^{-16}$) across a 106-point grid spanning $(s_1, s_2, R) \in [0.05, 2.0]^2 \times [1, 50]$ plus extreme asymptotic regimes; runtime $< 5\,\mathrm{s}$ on a Mac mini.

\bibitem{lemma-ct-higher-order} V.\ Bonfioli, Sage symbolic computation of the next-order Schwinger expansion for Lemma~\ref{lem:cross-term}: \texttt{verification/lemma\_ct\_higher\_order.sage}. Verifies $F(u, 0) = 1/(1+u^2)$ exactly, the symmetric series $F(\epsilon, \epsilon) = 1 - 2\epsilon^2 + 8(1+\log\epsilon)\epsilon^4 + O(\epsilon^6\log\epsilon)$, and the identification of the $\log(s/R)$ remainder at the $u^2 v^2$ cross-term.

\bibitem{lemma-ct-elementary} V.\ Bonfioli, symbolic and high-precision verification of the elementary and hyperbolic-distance closed forms \eqref{eq:lemma-ct-elementary}--\eqref{eq:lemma-ct-hyperbolic}: \texttt{verification/lemma\_ct\_elementary.py}. Verifies the pointwise integrand identity and the numerator identity $N = \Sigma D$ symbolically (exact); the closed form against 50-digit adaptive quadrature of the integral in \eqref{eq:lemma-cross-term-exact} on a 60-point grid (max.\ rel.\ error $< 10^{-45}$); the full inner product against the Rust 4-D grid values of \cite{lemma-ct-rust} (machine precision); the limits $G(0^+) = 2/3$ and $G(d) \sim 2e^{-d}$; and the symmetric expansion coefficient $8(1 + \log\epsilon)$ of \cite{lemma-ct-higher-order} directly from the closed form.

\bibitem{Sahbani1997} J.\ Sahbani, \emph{The conjugate operator method for locally regular Hamiltonians}, J.\ Operator Theory \textbf{38} (1997) 297--322.

\bibitem{Stein1970} E.M.\ Stein, \emph{Singular Integrals and Differentiability Properties of Functions}, Princeton Math.\ Series \textbf{30}, Princeton Univ.\ Press (1970). Ch.\ V, Thm.\ 1 and Cor.\ 2, pp.\ 119--121, for the weighted Riesz-composition / Hardy--Littlewood--Sobolev convolution estimate used in Lemma~\ref{app:lem:green} (the composition estimate is Cor.\ 2).

\bibitem{ReedSimonII} M.\ Reed, B.\ Simon, \emph{Methods of Modern Mathematical Physics II: Fourier Analysis, Self-Adjointness}, Academic Press (1975). Theorem X.39 (Nelson's analytic vector theorem) for essential self-adjointness of the cutoff dilation generator $A_j$ on $C_c^\infty$.

\bibitem{ABMG1996} W.O.\ Amrein, A.\ Boutet de Monvel, V.\ Georgescu, \emph{$C_0$-Groups, Commutator Methods and Spectral Theory of $N$-Body Hamiltonians}, Progress in Mathematics \textbf{135}, Birkh\"auser (1996).

\bibitem{FroeseHislop1989} R.\ Froese, P.\ Hislop, \emph{Spectral analysis of second-order elliptic operators on noncompact manifolds}, Duke Math.\ J.\ \textbf{58} (1989) 103--129.

\bibitem{mourre-cusp-script} V.\ Bonfioli, Sage symbolic verification of the $\bbH^5_r$ Mourre commutator identity and the iterated-commutator bound for the perturbation tensor: \texttt{verification/mourre\_estimate\_cusp.sage} and \texttt{verification/mourre\_iterated\_commutator.sage}.

\bibitem{lemma-od-script} V.\ Bonfioli, Sage derivation and numerical verification of the bubble-bubble cross-block closed forms for Lemma~\ref{lem:OD} (scale-position and the IR-structure of position-position): \texttt{verification/lemma\_od\_position\_terms.sage}. Computes the exact Schwinger closed form (with denominator $X(t)^2$ from the squared propagators after $\sigma$-integration), confirms the leading asymptotic \eqref{eq:OD-bare-sx-asymp} to relative error $<10^{-3}$ across $(R, s_1, s_2)$ ranging over four orders of magnitude, identifies the logarithmic infrared divergence of bare position-position cross-terms, and documents the gauge-projection argument leading to \eqref{eq:OD-horiz}.

\end{thebibliography}

\end{document}
